 %%%%%%%%%%%%%%%%%%%%%%%%%%%%%%%%%%%%%%%%%%%%%%%%%%%%%%%%%%
  %% MARK: art: The Invariant 1-Form
  %%%%%%%%%%%%%%%%%%%%%%%%%%%%%%%%%%%%%%%%%%%%%%%%%%%%%%%%%%
  
  \begin{article}\artlabel[The Prequantum $\pmb{1}$-Form $\pmb{\blambda}$ and its Curvature]
    \label{art: The Invariant 1-Form}
    
    By definition,
    the space of morphisms $\cY$ of the groupoid $\TT_\omega$ is the quotient of $\Paths(\X)$ by the equivalence relation $\sim_\omega$:
    $$
      \gamma \sim_\omega \gamma' \quad \text{iff} \quad
      \begin{cases}
      \ends(\gamma) = \ends(\gamma'), \\
      \phi(\gamma,\gamma') = 0,
      \end{cases}
      $$
    where $\phi : \ends^*(\Paths(\X)) \to \T_\omega$ is the cocycle defined in Article \ref{art: The Cocycle phi}.
    The condition $\phi(\gamma,\gamma') = 0$ is equivalent to
    $$
      \gamma \sim_\omega \gamma' \widetext{iff} \int_{\gamma'}^{\gamma} \Komega_{x,x'} \in \P_\omega.
      $$
    where $(x,x') = \ends(\gamma) = \ends(\gamma')$.
    
    \begin{theorem}{}
      There exists a unique differential $1$-form $\blambda$ on the space of morphisms $\cY = \Mor(\TT_\omega)$ such that $\class_\omega^*(\blambda) = \Komega$.
      We call $\blambda$ the \emph{prequantum $1$-form} on the groupoid $\TT_\omega$.
      Furthermore, $\blambda$ is invariant under left and right composition in $\TT_\omega$,
      and its curvature is given by $\der\blambda = \1^*\omega - \0^*\omega$,
      where $\1, \0 : \cY \to \X$ are the target and source maps.
    \end{theorem}
    
  \end{article}
  
  \begin{proof}
    \textbf{Existence and Uniqueness of $\blambda$.}~
    We use the criterion for a differential form on a space $\W$ to be the pullback of a differential form on a quotient space $\Z = \W/\!\sim$ via the projection $p : \W \to \Z$.
    A differential form $\alpha \in \Omega^k(\W)$ is the pullback of a unique differential form $\beta \in \Omega^k(\Z)$ (i.e., $\alpha = p^*(\beta)$) if and only if $\pr_1^*(\alpha) - \pr_2^*(\alpha)$ vanishes on the total space of the tautological pullback $p^*(\W) = \{ (w_1, w_2) \in \W \times \W \mid p(w_1) = p(w_2) \}$,
    where $\pr_1, \pr_2 : p^*(\W) \to \W$ are the projection maps \cite[\textsection 6.38 (Note 2)]{PIZ13}.
    
    In our case,
    $\W = \Paths(\X)$, $\Z = \cY$,
    $p = \class_\omega$,
    and $\alpha = \Komega \in \Omega^1(\Paths(\X))$.
    The total space of the tautological pullback is $\class_\omega^*(\Paths(\X)) = \{ (\gamma, \gamma') \in \Paths(\X)^2 \mid \class_\omega(\gamma) = \class_\omega(\gamma') \}$.
    By definition of $\sim_\omega$, $\class_\omega(\gamma) = \class_\omega(\gamma')$ if and only if $\gamma \sim_\omega \gamma'$,
    which means $\ends(\gamma) = \ends(\gamma')$ and $\phi(\gamma,\gamma') = 0$.
    Thus,
    $\class_\omega^*(\Paths(\X)) = \{ (\gamma, \gamma') \in \ends^*(\Paths(\X)) \mid \phi(\gamma,\gamma') = 0 \}$.
    
    We need to check if $\pr_1^*(\Komega) - \pr_2^*(\Komega)$ vanishes on this subspace.
    From Article \ref{art: The Cocycle phi},
    we know that $\pr_1^*(\Komega) - \pr_2^*(\Komega) = \phi^*(\theta)$ on $\ends^*(\Paths(\X))$,
    where $\theta$ is the canonical $1$-form on $\T_\omega$.
    Thus:
    $$
      (\pr_1^*(\Komega) - \pr_2^*(\Komega)) \restriction \class_\omega^*(\Paths(\X)) = (\phi^*(\theta)) \restriction \phi^{-1}(0).
      $$
    Since $\theta$ is a differential form on $\T_\omega$,
    its pullback $\phi^*(\theta)$ restricted to the inverse image of a point,
    that is,
    $(\phi^*(\theta)) \restriction \phi^{-1}(0)$,
    is zero because differential forms vanish on constant plots \cite[Ex. 96]{PIZ13}.
    Indeed,
    let $\P$ be a plot in $\phi^{-1}(0)$,
    that is $\phi \circ \P = 0$.
    Then,
    $\phi^*(\theta)(\P) = \theta(\phi \circ \P)$,
    but $\phi \circ \P(r) = 0$,
    thus $\phi^*(\theta)(\P) = 0$ for all $\P$ in $\phi^{-1}(0)$.
    The criterion is satisfied,
    which guarantees the existence and uniqueness of a differential $1$-form $\blambda \in \Omega^1(\cY)$ such that $\class_\omega^*(\blambda) = \Komega$.
    
    \textbf{Invariance of $\blambda$.}~
    We need to show that for any $y_0 \in \cY$,
    the pullback of $\blambda$ by left and right composition with $y_0$ is related to $\blambda$.
    Let $y_0 = [\gamma_0]_\omega \in \cY$.
    The left composition map $\eL(y_0) : \cY_{\trg(y_0), \ast} \to \cY_{\src(y_0), \ast}$ is given by $\eL(y_0)(y) = y_0 \cdot y$.
    Similarly,
    the right composition map $\eR(y_0) : \cY_{\ast, \src(y_0)} \to \cY_{\ast, \trg(y_0)}$ is given by $\eR(y_0)(y) = y \cdot y_0$.
    
    The invariance of $\Komega$ under pre-concatenation and post-concatenation in $\Paths(\X)$ is a known property \cite[\textsection 6.85]{PIZ13}.
    For any path $\gamma_0 \in \Paths(\X)$,
    the map $\eL(\gamma_0) : \Paths(\X)_{\trg(\gamma_0), \ast} \to \Paths(\X)_{\src(\gamma_0), \ast}$ given by $\eL(\gamma_0)(\gamma) = \gamma_0 \vee \gamma$ satisfies $\eL(\gamma_0)^*(\Komega) = \Komega$.
    Similarly, $\eR(\gamma_0)^*(\Komega) = \Komega$.
    
    The composition maps in $\cY$ are defined by factorization of path concatenation: $\eL([\gamma_0]_\omega)([\gamma]_\omega) = [\gamma_0 \vee \gamma]_\omega = \class_\omega(\eL(\gamma_0)(\gamma))$.
    This gives the commutative diagram:
    \begin{center}
      \begin{tikzcd}[column sep=large, row sep=large, every label/.append style = {font = \small}]
        \Paths(\X)_{\trg(\gamma_0), \ast} \arrow[d, swap, "\class_\omega"] \arrow[r, "\eL(\gamma_0)"]  & \Paths(\X)_{\src(\gamma_0), \ast} \arrow[d,"\class_\omega"] \\
        \cY_{\trg(y_0), \ast} \arrow[r,swap, "\eL(y_0)"] & \cY_{\src(y_0), \ast}
      \end{tikzcd}
    \end{center}
    Taking the pullback by $\eL(y_0)$ and using $\class_\omega^*(\blambda) = \Komega$:
    \begin{multline*}
      \eL(y_0)^*(\blambda) = \eL(y_0)^*(\class_\omega^*)^{-1}(\Komega) = (\class_\omega \circ \eL(\gamma_0))^*(\blambda) \\
      = \eL(\gamma_0)^*(\class_\omega^*(\blambda)) = \eL(\gamma_0)^*(\Komega).
    \end{multline*}
    By the invariance of $\Komega$ under $\eL(\gamma_0)$, $\eL(\gamma_0)^*(\Komega) = \Komega$.
    Thus,
    $\eL(y_0)^*(\blambda) = \Komega$.
    Since $\class_\omega^*(\blambda) = \Komega$,
    we have $\eL(y_0)^*(\blambda) = \class_\omega^*(\blambda)$.
    Restricting to the domain $\cY_{\src(y_0), \ast}$ and using the fact that $\class_\omega$ is a subduction from $\Paths(\X)_{\src(\gamma_0), \ast}$ to $\cY_{\src(y_0), \ast}$,
    we can push the equality down to $\cY$:
    $$
      \eL(y_0)^*(\blambda) \restriction {\cY_{\src(y_0), \ast}} = \blambda \restriction {\cY_{\trg(y_0), \ast}}.
      $$
    This shows left invariance.
    The proof for right invariance is analogous,
    using the postcomposition map $\eR$.
    
    \textbf{Curvature Formula.}
    With obvious notations:
    $\ends_{\cY} \circ \class_\omega = \ends_{\Paths}$.
    Then,
    Then,
    taking the pullback of $\pr_2^*(\omega) - \pr_1^*(\omega)$ by both sides of this equality,
    we get $\class_\omega^*(\ends_{\cY}^*(\pr_2^*(\omega) - \pr_1^*(\omega)))
    = \ends_{\Paths}^*(\pr_2^*(\omega) - \pr_1^*(\omega))
    = \1_{\paths}^*(\omega) - \0_{\paths}^*(\omega)
    = \der[\Komega] = \der[\class_\omega^*(\blambda)]
    = \class_\omega^*(\der\blambda)$.
    Since $\class_\omega$ is a subduction,
    we get $d\blambda = \ends_{\cY}^*(\pr_2^*\omega) - \ends_{\cY}^*(\pr_1^*\omega) = (\pr_2 \circ \ends_{\cY})^*\omega - (\pr_1 \circ \ends_{\cY})^*\omega = \1_{\cY}^*\omega - \0_{\cY}^*\omega$.
    which can also be written as $d\blambda + \ends_{\cY}^*(\omega \ominus \omega) = 0$.
  \end{proof}
  
  %%%%%%%%%%%%%%%%%%%%%%%%%%%%%%%%%%%%%%%%%%%%%%%%%%%%%%%%%%
  %% MARK: art: The Isotropy
  %%%%%%%%%%%%%%%%%%%%%%%%%%%%%%%%%%%%%%%%%%%%%%%%%%%%%%%%%%
  
  \begin{article}\artlabel[The Isotropy $\pmb{\TT_{\omega,x}}$]
    \label{art: The Isotropy}
    
    By construction of the groupoid $\TT_\omega$,
    the isotropy at a point $x \in \X$ is the subspace
    $$
      \TT_{\omega,x} := \Mor_{\TT_\omega}(x,x) = \big\{ [\ell]_\omega \mid \ell \in \Loops(\X,x) \big\}.
      $$
    Two loops $\ell$ and $\ell'$ based at $x$ are equivalent if and only if $\phi(\ell,\ell') = 0$.
    Define
    $$
      \eF: \Loops(\X,x) \to \T_\omega \widetext{by} \eF(\ell) = \pi_\omega \left(\int_{\bmx}^\ell \Komega \right),
      $$
    with $\bmx$ the constant loop $[t \mapsto x]$.
    
    \begin{theorem}{}
      Two loops $\ell,\ell' \in \Loops(\X,x)$ are equivalent if and only if $\eF(\ell) = \eF(\ell')$,
      and the map $\eF$ is a subduction.
      The projection $\overline\eF : \TT_{\omega,x} \to \T_\omega$,
      defined by $\overline\eF([\ell]_\omega) = \eF(\ell)$ is an isomorphism of diffeological groups.
      \begin{center}
        \begin{tikzcd}[column sep=large, row sep=large, every label/.append style = {font = \small}]
          \Loops(\X,x) \arrow[dr, "\eF"] \arrow[d, swap, "\class_\omega"] & {} \\
          \TT_{\omega,x} = \Loops(\X,x)/\!\sim_\omega \arrow[r, swap, "\overline\eF"] & \T_\omega
        \end{tikzcd}
      \end{center}
    \end{theorem}
  \end{article}
  
  \begin{proof}
    The notation $\TT_{\omega,x}$ refers to the isotropy group of the groupoid $\TT_\omega$ at the point $x \in \X$.
    That is,
    by definition:
    $$
      \TT_{\omega,x} = \{ [\ell]_\omega \mid \ell \in \Loops(\X,x) \}.
      $$
    Thus,
    $$
      \TT_{\omega,x} = \Loops(\X,x)/\!\sim_\omega,
      \widetext{with} \ell \sim_\omega \ell' \widetext{iff} \int_\ell^{\ell'} \Komega_{x,x} \in \P_\omega.
      $$
    Since $\pi_1(\X) = \{0\}$,
    $\ell$ and $\ell'$ are homotopic to the constant path $\bmx : t \mapsto x$.
    Moreover,
    by definition of $\P_\omega$,
    $$
      \int_{\bmx}^{\ell} \Komega_{x,x} + \int_{\ell}^{\ell'} \Komega_{x,x} + \int_{\ell'}^{\bmx} \Komega_{x,x} \in \P_\omega,
      $$
    since this sum of integrals corresponds to the integral of $\Komega_{x,x}$ on a loop in $\Loops(\X)$.
    Hence
    \begin{multline*}
      \int_{\ell}^{\ell'} \Komega_{x,x} \in \P_\omega
      \ \Leftrightarrow \
      \int_{\bmx}^{\ell'} \Komega_{x,x} - \int_{\bmx}^{\ell} \Komega_{x,x} \in \P_\omega \\
      \ \Leftrightarrow \
      \pi_\omega\left(\int_{\bmx}^{\ell'} \Komega_{x,x}\right) = \pi_\omega\left(\int_{\bmx}^{\ell} \Komega_{x,x}\right).
    \end{multline*}
    Remember that $\pi_\omega : \RR \to \T_\omega = \RR/\P_\omega$.
    Hence
    $$
      \ell \sim_\omega \ell' \ \Leftrightarrow \ \eF(\ell) = \eF(\ell')
      \widetext{with}
      \eF(\ell) = \pi_\omega\left(\int_{\bmx}^{\ell} \Komega_{x,x}\right).
      $$
    Hence,
    set theoretically,
    $\eF : \Loops(\X,x) \to \T_\omega$ identifies $\Val(\eF)$ with $\TT_{\omega,x}$:
    $$
      \TT_{\omega,x} \equiv \Val(\eF).
      $$
    \begin{lemma}{}
      The subset $\Val(\eF)$ is a subgroup of $\T_\omega$.
      It can only be $\{0\}$ or $\T_\omega$,
      but if $\omega \neq 0$,
      then $\Val(\eF) = \T_\omega$.
    \end{lemma}
    
    \textit{Proof of Lemma.}\hspace{1ex}
    Let $\Paths(\Loops(\X,x),\bmx,\ast)$ be the subspace of $\Paths(\Loops(\X))$ of homotopies connecting the constant loop $\bmx$ to any loop $\ell \in \Loops(\X,x)$.
    Define
    $$
      \FF : \Paths(\Loops(\X,x),\bmx,\ast) \to \RR, \text{ by } \FF(\sigma) = \int_\sigma \Komega,
      $$
    where $\sigma : s \mapsto \ell_s$ is a path in $\Loops(\X,x)$ connecting $\bmx$ to some loop $\ell$.
    Explicitly,
    $$
      \FF(\sigma) =  \int_0^1 \Komega(\sigma)_s(1) \ds = \int_0^1 \ds \int_0^1 \omega\left(\vect{t \\ s} \mapsto \ell_s(t)\right)_{t \choose s} \vect{1 \\ 0}\vect{0 \\ 1} \ dt
      $$
    Let $\sigma$ and $\sigma'$ be two elements of $\Paths(\Loops(\X,x),\bmx,\ast)$,
    connecting $\bmx$ to $\ell$ and $\bmx$ to $\ell'$.
    Let $\bar\sigma: s \mapsto \bar\ell_s = [t \mapsto \ell_s(1-t)]$,
    hence $\bar\sigma$ is a path in $\Loops(\X,x)$ connecting $\bmx$ to the reverse $\bar\ell$.
    By a change of variable $t \mapsto 1-t$ in the expression of $\FF(\sigma)$ above we have $\Komega(\bar\sigma) = - \Komega(\sigma)$.
    Also,
    since for all $s$ we have $\sigma(s)(1) = \sigma'(s)(0) = x$,
    we can consider the concatenation of homotopies $\sigma * \sigma'$.
    Then,
    thanks to Article~\ref{art: Integrating-concatenations-over-homotopies},
    and to the property of $\bar\sigma$,
    we have:
    \begin{equation}
      \renewcommand{\theequation}{$\ast$}
      \FF(\sigma*\sigma') = \FF(\sigma) + \FF(\sigma')
      \text{ and }
      \FF(\bar\sigma) = - \FF(\sigma).
    \end{equation}
    Thus,
    if $t = \FF(\sigma)$ and $t' = \FF(\sigma')$,
    then $t+t' = \FF(\sigma*\sigma') \in \Val(\FF)$,
    and $-t = \FF(\bar\sigma) \in \Val(\FF)$.
    Therefore,
    $\Val(\FF)$ is a subgroup of $\RR$.
    But,
    $\Val(\FF)$ is connected because the space of paths in $\Loops(\X)$ based at $\bmx$,
    i.e.,
    $\Paths(\Loops(\X,x),\bmx,\ast)$,
    is connected (actually contractible) and $\FF$ is smooth.
    Thus,
    either $\Val(\FF) = \{0\}$ or $\Val(\FF) = \RR$.
    Since $\Val(\eF) = \pi_\omega(\Val(\FF))$,
    then $\Val(\eF) = \{0\}$ or $\Val(\eF) = \T_\omega$.\hfill$\square$
    
    Let us examine the two different cases.
    
    \textbf{The zero case.}~If $\Val(\eF) = \{0\}$,
    then the isotropy $\TT_{\omega,x}$ is reduced to the identity $\{1_x\}$,
    and $\cY = \X \times \X$.
    Thus,
    $\der \blambda = \1^*(\omega) - \0^*(\omega)$.
    Restricted to $\cY_x = \{x\} \times \X$ we get $\omega = d\lambda$,
    and $\lambda = \blambda \restriction \{x\} \times \X$,
    that is,
    a $1$-form on $\X$.
    But returning to $\Komega$,
    we have $\Komega = \CHO(d \lambda) = \1^*(\lambda) -\0^*(\lambda) - d(\CHO\lambda)$,
    then  $\Komega \restriction \Loops(\X,x) = d(\CHO\lambda)$ and the condition $\eF = 0$ writes,
    for all $\ell \in \Loops(\X,x)$,
    that is:
    $$
      0 = \int_\bmx^\ell \Komega = \int_\bmx^\ell d[\CHO\lambda] = (\CHO\lambda)(\ell) - (\CHO\lambda)(\bmx) = \int_\ell \lambda - 0 = \int_\ell \lambda.
      $$
    Indeed,
    by definition of the chain-homotopy operator $\CHO$,
    for a $1$-form $\lambda$,
    $\CHO\lambda$ is the function from $\Paths(\X)$ to $\RR$:
    $$
      (\CHO\lambda)(\gamma) = \int_0^1 \lambda(\gamma)_t(1)\, dt = \int_\gamma \lambda.
      $$
    Therefore,
    $\lambda$ is a $1$-form vanishing on every loop,
    thus $\lambda$ is closed and $\omega = 0$;
    see \cite[Ex. 118]{PIZ13},
    a case we excluded.
    
    \textbf{The full case.}~If $\Val(\eF) = \T_\omega$.
    Then,
    $\omega \neq 0$ and $\eF : \Loops(\X,x) \to \T_\omega$ is a subduction projecting to a smooth isomorphism $\overline\eF : \TT_{\omega,x} \to \T_\omega$.
    
    This is what we will prove now.
    The map $\eF$ is already a smooth surjection,
    let us check that $\FF$ is a subduction,
    which will imply that $\eF$ itself is a subduction,
    and therefore that its projection $\overline{\eF} : \TT_{\omega,x} \to \T_\omega$ is a smooth isomorphism.
    
    Let us choose a path
    $$
      \sigma \in \Paths(\Loops(\X,x),\bmx,\ast)
      \widetext{such that} \FF(\sigma) \neq 0.
      $$
    Such path exists since $\FF: \Paths(\Loops(\X,x),\bmx,\ast) \to \RR$ is surjective.
    Let
    $$
      \sigma_s(t) = \sigma(st)
      \widetext{and} \varphi(s) = \FF(\sigma_s).
      $$
    So,
    $\varphi$ is a smooth real function such that
    $$
      \varphi(0) = 0
      \widetext{and}
      \varphi(1) = \FF(\sigma) \neq 0,
      $$
    since $\varphi(0) = \FF(\sigma_0)$ and $\sigma_0 = [t \mapsto \sigma(0)=\bmx]$.
    Thus,
    there exists
    $$
      s_0 \in \left]0,1\right[
      \widetext{such that}
      \varphi'(s_0) \neq 0,
      $$
    where $\varphi'$ denotes the derivative of $\varphi$.
    Otherwise $\varphi$ would be constant and then equal to $0$.
    Let
    $$
      \fS(s) = \bar\sigma_{s_0} * \sigma_{s+s_0},
      \widetext{where}
      \bar\sigma_{s_0}(t) = [ t' \mapsto \sigma_{s_0}(t)(1-t')].
      $$
    Define
    $$
      \psi = \FF \circ \fS.
      $$
    That is,
    $\psi \in \Cinfty(\RR,\RR)$ with:
    
    \begin{center}
      \begin{tikzcd}[column sep=small, row sep=large, every label/.append style = {font = \small}]
        {} & \Paths(\Loops(\X,x),\bmx,\ast) \arrow[dr, "\FF"] & {}\\
        % {} & \Paths(\Loops(\X,x),\bmx,\ast) \arrow[dr, start anchor = south east, "\FF"] & {}\\
        %\RR \arrow[ur,  end anchor=south west, "\fS"] \arrow[rr, swap, "\psi"] & {} & \RR
        \RR \arrow[ur, "\fS"] \arrow[rr, swap, "\psi"] & {} & \RR
      \end{tikzcd}
    \end{center}
    
    By Article \ref{art: Integrating-concatenations-over-homotopies}, we have:
    $$
      \FF(\bar\sigma_{s_0} * \sigma_{s+s_0}) = \FF(\sigma_{s+s_0}) - \FF(\sigma_{s_0}),
      $$
    and then
    $$
      \psi(s) = \FF(\fS(s)) = \FF(\bar\sigma_{s_0} * \sigma_{s+s_0}) = \FF(\sigma_{s+s_0}) - \FF(\sigma_{s_0}) = \varphi(s + s_0) - \varphi(s_0).
      $$
    Thus,
    $$
      \psi(0) = 0
      \widetext{and}
      \psi'(0) = \varphi'(s_0) \neq 0.
      $$
    Therefore,
    by the inverse function theorem on $\RR$,
    there exists $\varepsilon > 0$ such that $\psi \restriction \left]- \varepsilon,+ \varepsilon \right[$ is a local diffeomorphism mapping $\left]- \varepsilon,+ \varepsilon \right[$ to some open interval $\left]-a,+ b \right[$ with $a,b>0$ and $\psi^{-1}(0) = 0$.
    Hence,
    $$
      \psi \restriction  \left]- \varepsilon,+ \varepsilon \right[ = \FF \circ \fS \restriction  \left]- \varepsilon,+ \varepsilon \right[
      \widetext{implies}
      \psi \circ \psi^{-1} = \id_{\left]- \varepsilon,+ \varepsilon \right[} = \FF \circ (\fS \circ \psi^{-1}).
      $$
    Therefore,
    $$
      \FF \circ (\fS \circ \psi^{-1}) =\id_{\left]- \varepsilon,+ \varepsilon \right[}.
      $$
    Let us define:
    $$
      \Sigma = \fS \circ \psi^{-1},
      $$
    Hence,
    $$
      \begin{cases}
      \Sigma  : \left]-a,+ b \right[ \to \Paths(\Loops(\X,x),\bmx,\ast), \\
      \widetext{with}
      \Sigma(0) = \fS(0) = \bar\sigma_{s_0} \vee \sigma_{s_0}.
      \end{cases}
      $$
    is a local smooth section of $\FF : \Paths(\Loops(\X,x),\bmx,\ast) \to \RR$.
    $$
      \FF \circ \Sigma = \id_{\left]- \varepsilon,+ \varepsilon \right[},
      $$
    defined on an open interval around $0 \in \RR$ and mapping $0$ to $\fS(0) = \bar\sigma_{s_0} \vee \sigma_{s_0}$,
    and for any plot $\P$ in $\RR$ mapping $r_0$ to $0$,
    $\Q = \Sigma \circ \P$ is a lift in $\Paths(\Loops(\X,x),\bmx,\ast)$:
    \begin{center}
      \begin{tikzcd}[column sep=normal, row sep=large, every label/.append style = {font = \small}]
        \Paths(\Loops(\X,x),\bmx,\ast) \arrow[d, "\FF"] \\
        \left]-\epsilon,+ \epsilon \right[ \subset \RR  \arrow[u, bend left=30,"\Sigma = \fS \circ \psi^{-1}"]
      \end{tikzcd}
      \quad \text{implies} \quad
      \begin{tikzcd}[column sep=small, row sep=large, every label/.append style = {font = \small}]
        {} & \Paths(\Loops(\X,x),\bmx,\ast) \arrow[d, "\FF"] \\
        \U \arrow[r, swap, "\P"] \arrow[ur, "\Q"] & \left]-\epsilon,+ \epsilon \right[ \subset \RR  \arrow[u, bend left=30,"\Sigma"]
      \end{tikzcd}
    \end{center}
    Now,
    thanks to the additivity of function $\FF$,
    we can translate this local section everywhere on $\RR$.
    Let $t_0 \in \RR$,
    and recall that $\FF$ is surjective.
    Thus,
    there exists $\sigma' \in \Paths(\Loops(\X,x),\bmx,\ast)$ such that $\FF(\sigma') = t_0$.
    Let us then consider
    $$
      \fS'(s) = \sigma' \ast \fS(s),
      $$
    and let,
    $$
      \Psi(s) = \FF(\fS'(s)),
      \widetext{that is,}
      \FF(\sigma' \ast \fS(s)) = \FF(\sigma') + \FF(\fS(s)) = t_0 + \psi(s).
      $$
    So,
    $$
      \Psi(0) = t_0
      \widetext{and}
      \Psi'(0) = \psi'(0) \neq 0.
      $$
    Thus,
    $\Psi$ is a local diffeomorphism of $\RR$ mapping $0$ to $t_0$,
    its inverse $\Psi^{-1}$ maps $t_0$ to $0$.
    We define then:
    $$
      \Sigma' = \fS' \circ \Psi^{-1}
      $$
    Then,
    $\Sigma'$ is a local smooth section of $\FF$ over the interval $\left]t_0 - \epsilon, t_0 + \epsilon \right[$,
    $$
      \FF \circ \Sigma' = \id_{\left]t_0 - \epsilon, t_0 + \epsilon \right[},
      \widetext{such that}
      \Sigma'(t_0) = \sigma' \ast \fS(0) = \sigma' \ast (\bar\sigma_{s_0} \vee \sigma_{s_0}).
      $$
    Hence,
    we can lift locally $\FF$ around every point $t_0 \in \RR$.
    Therefore $\FF$ is a subduction.
    Then,
    by projection,
    $\eF = \pi_\omega \circ \FF$ is itself a subduction,
    and the map $\overline\eF : \TT_{\omega,x} \to \T_\omega$,
    which is injective by construction,
    is a diffeomorphism and therefore a smooth isomorphism from $\TT_{\omega,x}$ to $\T_\omega$.
    In conclusion,
    $\TT_\omega$ is a fibrating groupoid with isotropy $\T_\omega$.
  \end{proof}
  
  %%%%%%%%%%%%%%%%%%%%%%%%%%%%%%%%%%%%%%%%%%%%%%%%%%%%%%%%%%
  %%
  %% MARK: - sect: Symmetries of the Prequantum Groupoid
  %%
  %%%%%%%%%%%%%%%%%%%%%%%%%%%%%%%%%%%%%%%%%%%%%%%%%%%%%%%%%%
  
  %\section{Remarks and Applications}
  \section{\textsc{\textbf{Symmetries of the Prequantum Structure}}}
  \label{sec:Symmetries of the Prequantum Structure}
  
  In this section,
  we investigate the symmetries of the prequantum groupoid $(\TT_\omega, \blambda)$ constructed in the Main Theorem.
  We show that the group of automorphisms of this prequantum structure is naturally isomorphic to the group of $\omega$-preserving diffeomorphisms of the base space $\X$,
  demonstrating a faithful representation of the classical symmetries at the prequantum level.
  
  
  %%%%%%%%%%%%%%%%%%%%%%%%%%%%%%%%%%%%%%%%%%%%%%%%%%%%%%%%%%
  %% MARK: art: Automorphisms Induced by Symmetries
  %%%%%%%%%%%%%%%%%%%%%%%%%%%%%%%%%%%%%%%%%%%%%%%%%%%%%%%%%%
  
  \begin{article}
    \artlabel[Automorphisms Induced by Symmetries]
    \label{art:Automorphisms Induced by Symmetries}
    
    This article discusses the action of symmetries.
    The group of diffeomorphisms of $\X$ that preserve the closed $2$-form $\omega$,
    denoted by $\Diff(\X, \omega)$,
    plays a crucial role in the symmetry analysis of the system $(\X, \omega)$.
    In this framework,
    these symmetries lift naturally to automorphisms of the prequantum groupoid $(\TT_\omega, \blambda)$.
    Let us denote
    $$
      \Aut(\TT_\omega, \blambda) = \{ (\Phi,\phi) \in \Aut(\TT_\omega) \mid \Phi^*(\blambda) = \blambda \}.
      $$
    On the other hand,
    the group of automorphisms of $(\X,\omega)$ is denoted by
    $$
      \Diff(\X,\omega) = \{ \phi \in \Diff(\X) \mid \phi^*(\omega) = \omega \}.
      $$
    Note first that:
    
    \begin{proposition}{~1}
      If $(\Phi,\phi) \in \Aut(\TT_\omega,\blambda)$ then $\phi \in \Diff(\X,\omega)$.
    \end{proposition}
    
    And then,
    the construction of automorphisms of the prequantum groupoid:
    
    \begin{proposition}{~2.}
      Let $\TT_\omega$ be the associated prequantum groupoid and $\blambda$ its prequantum $1$-form,
      associated with the system $(\X,\omega)$,
      where $\X$ is connected and simply connected.
      
      For any $\phi \in \Diff(\X, \omega)$,
      that is,
      $\phi \in \Diff(\X)$ and $\phi^*(\omega) = \omega$,
      define the map
      $$
        \Phi: \cY \to \cY
        \quad \text{with} \quad
        \Phi([\gamma]_\omega) = [\phi \circ \gamma]_\omega.
        $$
      Then,
      this automorphism preserves the prequantum $1$-form $\blambda$,
      $\Phi^* \blambda = \blambda$,
      and
      $$
        (\Phi, \phi) \in \Aut(\TT_\omega,\blambda).
        $$
    \end{proposition}
    
    Note that this proposition shows that the map $(\Phi,\phi) \mapsto \phi$ from $\Aut(\TT_\omega,\blambda)$ to $\Diff(\X,\omega)$ is surjective.
    
  \end{article}
  
  \begin{proof}
    We first show that $\Phi$ is well-defined.
    Suppose $\gamma \sim_\omega \gamma'$.
    By definition,
    let $(x,x') = \ends(\gamma) = \ends(\gamma')$.
    Then,
    for any fixed-ends homotopy $\sigma : s \mapsto \gamma_s$ in $\ends^{-1}(x,x')$,
    from $\gamma = \gamma_0$ to $\gamma' = \gamma_1$,
    $$
      \int_\sigma \Komega \in \P_\omega.
      $$
    Consider the path $\phi \circ \sigma: \RR \to \Paths(\X)$ defined by $s \mapsto \phi \circ \gamma_s$.
    This is a fixed-ends homotopy from $\phi \circ \gamma$ to $\phi \circ \gamma'$,
    with $\phi\circ\gamma_s(0) = \phi(x)$ and $\phi\circ\gamma_s(1) = \phi(x')$.
    That is,
    $\phi \circ \sigma$ is a path in $\ends^{-1}(\phi(x),\phi(x'))$.
    Therefore,
    $$
      \int_{\phi \circ \sigma} \Komega = \int_\sigma \phi^*(\Komega) = \int_\sigma \Komega \in \P_\omega,
      $$
    thanks to the variance of the chain-homotopy operator \cite[\textsection 6.84]{PIZ13}.
    Thus,
    $\Phi$ is well-defined.
    
    Next,
    $\Phi$ is a smooth map from $\cY$ to $\cY$.
    Indeed,
    the map $\Psi: \Paths(\X) \to \Paths(\X)$ defined by $\Psi(\gamma) = \phi \circ \gamma$ is smooth,
    as it is the composition operator in the functional diffeology.
    Since $\class_\omega: \Paths(\X) \to \cY$ is a subduction and $\Phi \circ \class_\omega = \class_\omega \circ \Psi$,
    $\Phi$ is smooth because $\Psi$ is smooth.
    
    The map $\Phi$ is bijective because $\phi$ is a diffeomorphism.
    Its inverse is given by $\Phi^{-1}([\gamma']_\omega) = [\phi^{-1} \circ \gamma']_\omega$.
    
    Finally,
    we verify the functor properties:
    
    \textbf{(a)} Preservation of Source and Target:
    For $[\gamma]_\omega \in \cY$,
    \begin{multline*}
      \ends \circ \Phi([\gamma]_\omega)
      = \ends([\phi \circ \gamma]_\omega)
      = \ends(\phi \circ \gamma) \\
      = (\phi(\gamma(0)), \phi(\gamma(1)))
      = (\phi \times \phi)(\ends[\gamma]_\omega).
    \end{multline*}
    
    \textbf{(b)} Preservation of Composition:
    For composable $[\gamma]_\omega, [\gamma']_\omega \in \cY$,
    \begin{multline*}
      \Phi([\gamma]_\omega \cdot [\gamma']_\omega)
      = \Phi([\gamma \vee \gamma']_\omega)
      = [\phi \circ (\gamma \vee \gamma')]_\omega \\
      = [(\phi \circ \gamma) \vee (\phi \circ \gamma')]_\omega
      = [\phi \circ \gamma]_\omega \cdot [\phi \circ \gamma']_\omega
      = \Phi([\gamma]_\omega) \cdot \Phi([\gamma']_\omega).
    \end{multline*}
    
    \textbf{(c)} Preservation of Identity:
    For $x \in \X$,
    with $\bmx = [t \mapsto x]$,
    $$
      \Phi(\id_x) = \Phi([\bmx]_\omega) = [\phi \circ \bmx]_\omega = [\widehat{\phi(x)}]_\omega = \id_{\phi(x)}.
      $$
    
    Since $(\phi, \Phi)$ is a groupoid morphism and both $\phi$ and $\Phi$ are diffeomorphisms,
    $(\phi, \Phi)$ is a diffeological groupoid automorphism of $\TT_\omega$. This proves part (1).
    
    For part (2),
    we show that $\Phi^* \blambda = \blambda$.
    Let $\phi_* : \Paths(\X) \to \Paths(\X)$ defined by $\phi_*(\gamma) = \phi \circ \gamma$.
    The map $\Phi$ is defined by the commutative diagram:
    \begin{center}
      \begin{tikzcd}[column sep=large, row sep=large, every label/.append style = {font = \small}]
        \Paths(\X) \arrow[d, swap, "\class_\omega"] \arrow[r, "\phi_*"]  & \Paths(\X) \arrow[d,"\class_\omega"] \\
        \cY \arrow[r,swap, "\Phi"] & \cY
      \end{tikzcd}
    \end{center}
    Then,
    $(\Phi \circ \class_\omega)^*(\blambda) = (\class_\omega \circ \phi_*)^*(\blambda)$.
    That is,
    $\class_\omega^*(\Phi^*(\blambda)) = (\phi_*)^*(\class_\omega^*(\blambda))$.
    Thus,
    $\class_\omega^*(\Phi^*(\blambda)) = (\phi_*)^*(\Komega)$.
    But,
    since $\phi^*(\omega) = \omega$,
    thanks to the variance of the chain-homotopy operator $\CHO$,
    see \cite[\textsection 6.84]{PIZ13},
    $(\phi_*)^*(\Komega) = \Komega$,
    hence:
    $\class_\omega^*(\Phi^*(\blambda)) = \Komega = \class_\omega(\blambda)$.
    Thus,
    since $\class_\omega$ is a subduction,
    thanks to \cite[\textsection 6.39]{PIZ13},
    $\Phi^*(\blambda) = \blambda$.
  \end{proof}
  
  
  %%%%%%%%%%%%%%%%%%%%%%%%%%%%%%%%%%%%%%%%%%%%%%%%%%%%%%%%%%
  %% MARK: art: The Automorphism Group
  %%%%%%%%%%%%%%%%%%%%%%%%%%%%%%%%%%%%%%%%%%%%%%%%%%%%%%%%%%
  
  \begin{article}
    \artlabel[The Automorphism Group of\ $\pmb{(\TT_\omega,\lambda)}$]
    \label{art: The Automorphism Group}
    
    In this article,
    we discuss the automorphism group of the structure $(\TT_\omega,\blambda)$ and show that its group of automorphisms is isomorphic to the group of symmetries of the structure $(\X,\omega)$,
    i.e.,
    $\Diff(\X,\omega)$.
    
    The group of automorphisms of the groupoid $\TT_\omega$ preserving the prequantum $1$-form $\blambda$ is denoted by $\Aut(\TT_\omega,\blambda)$.
    An element $(\Phi,\phi) \in \Aut(\TT_\omega,\blambda)$ is an element $(\Phi,\phi) \in \Aut(\TT_\omega)$ such that $\Phi^*(\blambda) = \blambda$.
    That is:
    %
    \begin{enumerate}
      \item $\Phi \in \Diff(\cY)$,
      with $\cY = \Mor(\TT_\omega)$.
      \item There exists $\phi \in \Diff(\X)$ such that: $\ends \circ\; \Phi = (\phi \times \phi) \circ \ends$,
      
      \begin{tikzcd}[column sep=large, row sep=large, every label/.append style = {font = \small}]
        \cY \arrow[d, swap, "\ends"] \arrow[r, "\Phi"]  & \cY \arrow[d,"\ends"] \\
        \X \times \X \arrow[r,swap, "\phi \times \phi"] & \X \times \X
      \end{tikzcd}
      \quad
      $\begin{array}{c}
      \text{with} \\
      (\phi \times \phi)(x,x') = (\phi(x),\phi(x')).
      \end{array}$
      
      \item $\Phi(y \cdot y') = \Phi(y) \cdot \Phi(y')$,
      for all juxtaposable pairs.
      \item $\Phi^*(\blambda) = \blambda$.
    \end{enumerate}
    
    \begin{proposition}{}
      Let $(\Phi,\phi)$ be a pair in $\Aut(\TT_\omega)$.
      If $\Phi^*(\blambda) = \blambda$,
      then $\phi \in \Diff(\X,\omega)$.
    \end{proposition}
    
    We finally establish the most important result on the group of automorphisms of the prequantum groupoid:
    
    \begin{theorem}{}
      The projection $(\Phi,\phi) \mapsto \phi$ from the group of automorphisms of the prequantum groupoid $\Aut(\TT_\omega,\blambda)$ to the group of symmetries of the parasymplectic form $\omega$ on $\X$,
      $\Diff(\X,\omega)$,
      is an isomorphism of diffeological groups:
      $$
        \Aut(\TT_\omega, \blambda) \simeq \Diff(\X,\omega).
        $$
    \end{theorem}
    
    \textbf{Note.} This theorem reveals a particularly interesting situation:
    the group of $\omega$-preserving diffeomorphisms of the base space $(\X, \omega)$ has a full and faithful representation as automorphisms of the prequantum groupoid $(\TT_\omega, \blambda)$.
    This is a significant feature,
    as it implies that every symmetry of the classical system $(\X, \omega)$ lifts to a unique automorphism of the prequantum structure,
    and the entire group $\Diff(\X, \omega)$ is represented at the prequantum level without any loss of information or the need for a central extension.
    The prequantum property of the system $(\X,\omega)$ is captured in the isotropy groups $\TT_{\omega,x}$ of the groupoid $\TT_\omega$.
    This is achieved without abandoning any symmetries of the structure.
    
  \end{article}
  
  \begin{proof}
    Let us begin by proving the first assertion.
    Let us apply $\ends \circ \Phi = (\phi \times \phi) \circ \ends$ to the pullback of $\omega \ominus \omega$\,:
    \begin{align*}
      (\ends \circ\; \Phi)^*(\omega \ominus \omega) & = ((\phi \times \phi) \circ \ends)^*(\omega \ominus \omega) \\
      \Phi^*(\ends^*(\omega \ominus \omega)) & = \ends^*((\phi \times \phi)^*(\omega \ominus \omega)) \\
      \Phi^*(\der\blambda) & = \ends^*(\phi^*(\omega) \ominus \phi^*(\omega)) \\
      \der[\Phi^*\blambda] & = \ends^*(\phi^*(\omega) \ominus \phi^*(\omega)) \\
      \der\blambda & = \ends^*(\phi^*(\omega) \ominus \phi^*(\omega)) \\
      \ends^*(\omega \ominus \omega) & = \ends^*(\phi^*(\omega) \ominus \phi^*(\omega))
    \end{align*}
    Since $\ends$ is a subduction,
    this implies $\phi^*(\omega) \ominus \phi^*(\omega) = \omega \ominus \omega$,
    see \cite[\textsection 6.39]{PIZ13}.
    Now apply that to a plot $\P \times \Const$,
    where $\Const$ is any constant plot.
    Then we get $\phi^*(\omega)(\P) = \omega(\P)$ for all plots in $\X$,
    i.e.,
    $\phi^*(\omega) = \omega$.
    
    Now,
    let us consider the projection $\pr_2 (\Phi,\phi) \mapsto \phi$ from $\Aut(\TT_\omega,\blambda)$ to $\Diff(\X,\omega)$.
    We know already by Article \ref{art:Automorphisms Induced by Symmetries} that this is a surjection.
    Let us consider the kernel $\ker(\pr_2)$ of this projection,
    that is the subgroup of $(\Phi,\id_\X) \in \Aut(\TT_\omega,\blambda)$.
    Since $\Phi$ projects onto the identity of $\X$,
    one has:
    $$
      \ends \circ\ \Phi = \ends.
      $$
    That is,
    for all $y \in \cY$,
    $\ends(\Phi(y)) = \ends(y)$.
    Thus,
    $\Phi(y)$ is composable with $y^{-1}$,
    and $\tau(y) = \Phi(y) \cdot y^{-1}$ belongs to the isotropy group $\TT_{\omega,x}$,
    with $x = \src(y)$.
    By construction,
    the map $\tau$ is smooth.
    Therefore the map $\Phi \in \ker(\pr_2)$ writes $\Phi(y) = \tau(y) \cdot y$,
    and satisfies $\Phi^*(\blambda) = \blambda$.
    
    Now choose $x \in \X$,
    and let $\1_x : \Y_x \to \X$,
    where $\Y_x = \Mor_{\TT_\omega}(x,\ast)$ and $\1_x = \1$.
    Let $\lambda_x = \blambda \restriction \Y_x$.
    The projection $\1_x$ is a principal fibration with group $\T_x = \TT_{\omega,x}$,
    and the $1$-form $\lambda_x$ is actually a diffeological connection form satisfying the conditions \cite[\textsection 8.37]{PIZ13}:
    \begin{itemize}
      \item[$\ast$] The form is \emph{invariant} by $\T_x$, $\tau^*(\lambda_x) = \lambda_x$,
      for all $\tau \in \T_x$.
      \item[$\ast$] The form is \emph{calibrated}:
      $\hat y^*(\lambda_x) = \theta$,
      where $\hat y : \tau \mapsto \tau(y)$ is the \emph{orbit map},
      and $\theta$ is the canonical $1$-form on $\T_x \simeq \RR/\P_\omega$.
    \end{itemize}
    Now, the restriction $\Phi_x$ of $\Phi$ on $\Y_x$ writes the same way,
    that is,
    $\Phi_x(y) = \tau(y) \cdot y$.
    
    \begin{lemma}{}
      The automorphism $\Phi_x : y \mapsto \tau(y) \cdot y$ preserves $\lambda_x$,
      i.e.,
      $\Phi_x^*(\lambda_x) = \lambda_x$,
      if and only if $\tau(y) = \tau_x$ is constant on $\Y_x$.
    \end{lemma}
    
    \smallskip
    \textit{Proof of the lemma.}\hspace{1ex}
    Let us compute $\Phi_x^*(\lambda_x)$ on a plot $r \mapsto y_r$.
    That is,
    $$
      \Phi_x^*(\lambda_x)(r \mapsto y_r)_r(\delta r) = \lambda_x(r \mapsto \tau(y_r) \cdot y_r)_r(\delta r),
      $$
    for all $r$ in the domain of the plot and $\delta r$ a tangent vector at the point $r$.
    For that,
    we will separate the variables:
    for any integers $n$ and $m$,
    any $n$-plot $\btau : \U \to \T_x$ and any $m$-plot $\P : \V \to \Y_x$,
    let $\btau \cdot \P$ be the plot of $\Y_x$ defined by
    $$
      \btau \cdot \P : \U \times \V \to \Y_x,  \text{ with }  \btau \cdot \P : (r,s) \mapsto \btau(r) \cdot \P(s).
      $$
    Let $(r,s) \in \U \times \V$ and two vectors $\delta r \in \RR^n$ and  $\delta s \in \RR^m$.
    The value of $\lambda_x$ on $\btau \cdot \P$ is given by
    $$
      \lambda_x(\btau \cdot \P)_{r \choose s} \vect{\delta r \\ \delta s} = \btau^*(\theta)_r(\delta r) + \lambda_x(\P)_s(\delta s).
      $$
    Indeed,
    let us develop the $1$-form $\lambda_x(\btau \cdot \P)$ defined on $\U \times \V$,
    \begin{eqnarray*}
      \lambda_x(\btau \cdot \P)_{r \choose s} \vect{\delta r \\ \delta s} & = & \lambda_x[(r,s) \mapsto (\btau(r),\P(s)) \mapsto \tau(r) \cdot \P(s)]_{r \choose s}\vect{\delta r \\ \delta s} \\
      & = & [\lambda_x^{\U,s}(r) \quad \lambda_x^{\V,r}(s)] \vect{\delta r \\ \delta s},
    \end{eqnarray*}
    because every $1$-form on $\U \times \V$ at a point $(r,s) \in \U \times \V$ writes $[\lambda_x^{\U,s}(r) \ \lambda_x^{\V,r}(s)]$,
    where $\lambda_x^{\U,s}$ is a $1$-form on $\U$ depending on $s$,
    and $\lambda_x^{\V,r}$ is a $1$-form on $\V$ depending on $r$.
    Let us use indifferently $\tau_r = \btau(r)$,
    and $y_s = \P(s)$.
    We have
    \begin{eqnarray*}
      \lambda_x(\btau \cdot \P)_{r \choose s} \vect{ \delta r \\ \delta s} & = & \lambda_x^{\U,s}(r) (\delta r) + \lambda_x^{\V,r}(s)(\delta s) \\
      & = & \lambda_x[r \mapsto \btau(r) \cdot y_s]_r(\delta r) + \lambda_x[s \mapsto \tau_r \cdot \P(s)]_s(\delta s) \\
      & = & \lambda_x(\hat y_s \circ \btau)_r(\delta r) + \tau_r^*(\lambda)(\P)_s(\delta s) \\
      & = & \btau^*[{\hat y}_s^*(\lambda)]_r(\delta r) + \lambda(\P)_s(\delta s) \\
      & = & \btau^*(\theta)_r(\delta r) + \lambda(\P)_s(\delta s).
    \end{eqnarray*}
    Therefore,
    $$
      \Phi_x^*(\lambda_x)(r \mapsto y_r)_r(\delta r) = \btau^*(\theta)_r(\delta r) + \lambda(r \mapsto y_r)_r(\delta r),
      $$
    where $\btau(r) = \tau(y_r)$.
    Now,
    having $\Phi_x^*(\lambda_x) = \lambda_x$ means that $\Phi_x^*(\lambda_x)(r \mapsto y_r)_r(\delta r) = \lambda_x(r \mapsto y_r)_r(\delta r)$,
    for all plots $r \mapsto y_r$.
    That is,
    $\btau^*(\theta)_r(\delta r) = 0$ for all plots $r \mapsto y_r$.
    Locally,
    the plot $r \mapsto \btau(r) = \tau(y_r)$ lifts to $\RR$ into a plot $r \mapsto t(y_r)$,
    such that $\tau(y_r) = \class(t(y_r))$,
    where $\class : \RR \to \T_x = \RR/\P_\omega$.
    Then,
    $\btau^*(\theta)_r = [r \mapsto t(y_r)]^*(\dt) = \der[r \mapsto(t(y_r))]$.
    And this must be zero.
    Thus $r \mapsto t(y_r)$ is locally constant,
    and so is $r \mapsto \tau(y_r)$,
    for all plots $r \mapsto y_r$.
    Hence,
    $y \mapsto \tau(y)$ is locally constant,
    and since $\Y_x$ is connected,
    $y \mapsto \tau(y)$ is constant,
    equal to some $\tau_x$.\hfill $\square$
    
    Now,
    the morphism $\Phi$ writes $\Phi(y) = \tau_x \cdot y$,
    with $x = \src(y)$.
    Applying the groupoid morphism rule $\Phi(y \cdot y') = \Phi(y) \cdot \Phi(y')$,
    we get,
    with $x = \src(y)$ and $x' = \trg(y) = \src(y')$,
    $$
      \tau_x \cdot y \cdot y' = \tau_x \cdot y \cdot \tau_{x'} \cdot y' \quad \Rightarrow \quad \tau_{x'} = \id_{x'}.
      $$
    Therefore,
    for all $y$ and $x = \src(y)$,
    $\Phi(y) = y$.
    The only automorphism of the structure $(\TT_\omega,\blambda)$ projecting onto the identity of $\X$ is the identity of $\cY$,
    and thus the projection $\pr_2: \Aut(\TT_\omega,\blambda) \to \Diff(\X,\omega)$ is an isomorphism.
  \end{proof}
  
  
  %%%%%%%%%%%%%%%%%%%%%%%%%%%%%%%%%%%%%%%%%%%%%%%%%%%%%%%%%%
  %%
  %% MARK: - sect: Remarks and Applications
  %%
  %%%%%%%%%%%%%%%%%%%%%%%%%%%%%%%%%%%%%%%%%%%%%%%%%%%%%%%%%%
  
  %\section{Remarks and Applications}
  \section{\textsc{\textbf{Remarks and Applications}}}
  \label{sec:Remarks and Applications}
  This section discusses the significance and implications of the main theorem,
  structured into several articles.
  
  %%%%%%%%%%%%%%%%%%%%%%%%%%%%%%%%%%%%%%%%%%%%%%%%%%%%%%%%%%
  %% MARK: art: On the System $TT_\omega$
  %%%%%%%%%%%%%%%%%%%%%%%%%%%%%%%%%%%%%%%%%%%%%%%%%%%%%%%%%%
  
  \begin{article}
    \artlabel[On the System: The Prequantum Groupoid as the Central Object]
    \label{art: On the System}
    
    The construction presented in this paper introduces a shift in the geometric quantization program by positioning the prequantum groupoid $(\TT_\omega, \blambda)$ as the central object,
    in place of the traditional prequantum principal bundle.
    
    This groupoid object is an intrinsic construction,
    derived solely from the parasymplectic space $(\X, \omega)$,
    as a diffeological quotient of the space of paths,
    a feature that incidentally resonates with Feynman's path integral approach to quantization.
    This single object serves as a unified geometric container possessing several key virtues:
    
    $\ast$~\textbf{It contains the original classical system $(\X, \omega)$}:
    The space of objects $\Obj(\TT_\omega)$ is $\X$,
    and the fundamental relationship defined by the $2$-form $\omega$ is encoded in the curvature of the prequantum $1$-form $\blambda$ on the space of morphisms $\Mor(\TT_\omega)$.
    
    $\ast$~\textbf{It embodies the prequantum total space}:
    The space of morphisms $\cY = \Mor(\TT_\omega)$ carries the fundamental prequantum $1$-form $\blambda$,
    which serves as the potential for $\omega$ according to the formula $d\blambda + \ends^*(\omega \oplus \omega) = 0$.
    
    $\ast$~\textbf{It faithfully represents the full symmetry group $\pmb{\Diff(\X, \omega)}$}:
    As shown in Article \ref{art: The Automorphism Group},
    the group of $\omega$-preserving diffeomorphisms of $\X$ acts as the full group of automorphisms of the entire $(\TT_\omega, \blambda)$ structure in a faithful manner,
    without involving central extensions at this level.
    
    $\ast$~\textbf{It geometrically captures the quantum phase information}:
    The periodicity associated with $\omega$ is intrinsically represented by the structure of the isotropy groups $\TT_{\omega,x}$ at each point $x \in \X$
    ---~the ``vertical'' structure of the groupoid which is isomorphic to the torus of periods $\T_\omega$~---
    as the quotient of the space of loops.
    
  \end{article}
  
  %%%%%%%%%%%%%%%%%%%%%%%%%%%%%%%%%%%%%%%%%%%%%%%%%%%%%%%%%%
  %% MARK: art: Structure of the Space of Morphisms
  %%%%%%%%%%%%%%%%%%%%%%%%%%%%%%%%%%%%%%%%%%%%%%%%%%%%%%%%%%
  
  \begin{article}
    \artlabel[Structure of the Space of Morphisms $\pmb{\cY}$]
    \label{art: Structure of the Space of Morphisms}
    
    The space of morphisms $\cY = \Mor(\TT_\omega)$ possesses a rich internal structure related to the principal bundle slices $\cY_x = \Mor_{\TT_\omega}(x, \ast)$,
    where $x \in \X$ is a fixed base point.
    Recall from the main theorem (Theorem \ref{sec:Main Theorem: The Prequantum Groupoid}, part 1) that $\1_x: \cY_x \to \X$ is a principal $\TT_{\omega,x}$-bundle with structure group $\TT_{\omega,x} \simeq \T_\omega$.
    
    Now,
    consider the square of the projection $\1_x : \cY_x \to \X$,
    that is,
    $$
      \1_x \times \1_x : \cY_x \times \cY_x \to \X \times \X,
      \widetext{with}
      \1_x \times \1_x  : (y,y') \mapsto (\1(y),\1(y')).
      $$
    Since $\1_x : \cY_x \to \X$ is itself a $\TT_{\omega,x}$-principal fiber bundle,
    this is a $\TT_{\omega,x}^2$-principal fiber bunlde,
    for the action $(\tau,\tau') \cdot (y,y') = (\tau \cdot y, \tau' \cdot y')$.
    According to the reduction of principal bundle described in \cite[\textsection 8.18]{PIZ13},
    the quotient of $\cY_x \times \cY_x$ by the diagonal subgroup $\TT_{\omega,x} \simeq \{(\tau,\tau) \mid \tau \in \TT_{\omega,x} \}$ is itself a fiber bundle over the same base space $\X \times \X$,
    by the projection $\class: (y,y') \mapsto \cY \times_{\TT_{\omega,x}} \cY = (\cY \times \cY)/{\TT_{\omega,x}}$,
    with fiber ${\TT_{\omega,x}}^2/{\TT_{\omega,x}} \simeq {\TT_{\omega,x}}$ (as homogeneous space).
    
    Now,
    the quotient $\cY_x \times_{\TT_{\omega,x}} \cY_x$ is naturally realized by the projection
    $$
      \class:  \cY_x \times \cY_x \to \cY
      \widetext{with}
      \class(y,y') = y^{-1} \cdot y'.
      $$
    Therefore,
    the projection $\ends: \cY \to \X \times \X$ is itself a fiber bundle with homogeneous fiber $\T_\omega$,
    which is summarized by the diagram:
    
    \begin{center}
      \begin{tikzcd}[column sep=large, row sep=large, every label/.append style = {font = \small}]
        \cY_x \times \cY_x \arrow[dr, swap, "\1_x \times \1_x"] \arrow[rr, "\class"]  & {} & \cY \arrow[dl,"\ends"] \\
        {} & \X \times \X & {}
      \end{tikzcd}
    \end{center}
    
    \smallskip
    \textbf{Note.}
    As a consequence of the groupoid structure,
    for any fixed path $\gamma_0 \in \Paths(\X,x,x')$,
    the map $\tau \mapsto \tau \cdot [\gamma_0]_\omega$ establishes a diffeomorphism from the isotropy group $\TT_{\omega,x}$ to the fiber $\Mor_{\TT_\omega}(x,x')$.
    This means that $\Mor_{\TT_\omega}(x,x')$ is a principal homogeneous space for the isotropy group $\TT_{\omega,x} \simeq \T_\omega$.
    
    Thus,
    the fibers of the bundle $\ends: \cY \to \X \times \X$ are all diffeomorphic to the torus of periods $\T_\omega$.
    This holds for any pair of points $(x,x')$.
    
    When the torus of periods $\T_\omega = \RR/\P_\omega$ is a strict diffeological space (i.e., not diffeomorphic to $\S^1$ or $\RR$),
    the total space of the groupoid $\cY$ will also be a strict diffeological space,
    even if the base space $\X$ is a smooth manifold.
    This highlights how the diffeological framework naturally accommodates these structures that arise from the geometry of $\omega$ and may not fit into the traditional manifold setting.
    
  \end{article}
  
  %%%%%%%%%%%%%%%%%%%%%%%%%%%%%%%%%%%%%%%%%%%%%%%%%%%%%%%%%%
  %% MARK: art: Exact Case: A Unique Prequantum Structure
  %%%%%%%%%%%%%%%%%%%%%%%%%%%%%%%%%%%%%%%%%%%%%%%%%%%%%%%%%%
  
  \begin{article}
    \artlabel[The Exact Case]
    \label{art: Exact Case}
    
    When the closed $2$-form $\omega$ is exact,
    i.e., $\omega = d\alpha$ for some $1$-form $\alpha \in \Omega^1(\X)$,
    the unique prequantum groupoid construction simplifies significantly,
    revealing a particularly direct relationship to the primitive $\alpha$.
    
    In this case,
    the group of periods $\P_\omega$ is trivial,
    $\P_\omega = \{0\}$.
    Indeed,
    $\Komega = \CHO(d\alpha) = \1^*(\alpha) - \0^*(\alpha) - d[\Kalpha]$.
    Note that $f_\omega := \Kalpha \in \Omega^0(\Paths(\X))$ is a smooth map,
    explicitly:
    $$
      f_\omega : \Paths(\X) \to \RR
      \widetext{is defined by}
      f_\omega(\gamma) = \int_\gamma \alpha.
      $$
    Now,
    restricted to $\Loops(\X)$,
    we get $\Komega \restriction \Loops(\X) = df_\omega \restriction \Loops(\X)$,
    and the periods are the integrals
    $\int_\sigma df_\omega = f_\omega(\sigma(1)) - f_\omega(\sigma(0))$.
    However,
    when $\sigma \in \Loops(\Loops(\X))$,
    $f_\omega(\sigma(1)) = f_\omega(\sigma(0))$,
    which implies $\P_\omega = \{0\}$ and $\T_\omega = \RR/\P_\omega = \RR$.
    
    Next,
    the equivalence relation $\gamma \sim_\omega \gamma'$ means:
    $\ends(\gamma) = \ends(\gamma') = (x,x')$ and $\int_\sigma \Komega = 0$ for any fixed-ends homotopy $\sigma : s \mapsto \gamma_s$ from $\gamma$ to $\gamma'$.
    But:
    \begin{multline*}
      \int_\sigma \Komega = \int_\sigma (\1^*(\alpha) - \0^*(\alpha) - d[f_\omega])
      = \int_{\1 \circ \sigma} \alpha -  \int_{\0 \circ \sigma} \alpha - \int_\sigma df_\omega \\
      = \int_{s \mapsto \gamma_s(1)} \alpha -  \int_{s \mapsto \gamma_s(0)} \alpha - (f_\omega(\gamma) - f_\omega(\gamma'))
      = \underbrace{\int_{s \mapsto x'} \alpha}_{=0} - \underbrace{\int_{s \mapsto x} \alpha}_{=0}  - f_\omega(\gamma) + f_\omega(\gamma')\\
      = f_\omega(\gamma') - f_\omega(\gamma)
    \end{multline*}
    Then,
    $$
      \gamma \sim_\omega \gamma'
      \widetext{iff}
      \ends(\gamma) = \ends(\gamma')
      \widetext{and}
      \int_\gamma \alpha = \int_{\gamma'} \alpha.
      $$
    
    The space of morphisms $\cY = \Paths(\X)/\!\sim_\omega$ can be identified diffeomorphically with $\X \times \RR \times \X$,
    by the subduction
    $$
      \cY \simeq \X \times \RR \times \X
      \widetext{with}
      \class_\omega : \gamma \mapsto \left(x = \gamma(0), t = \int_\gamma \alpha,x' = \gamma(1) \right)
      $$
    By additivity of the integral of $\alpha$ on paths,
    the groupoid composition becomes
    $$
      (x,t,x') \cdot (x',t',x'') = (x, t + t', x")
      $$
    This groupoid is an \emph{additive groupoid} over $\X$.
    Note that,
    the identity of the object $x \in \X$ in the groupoid is the arrow $(x,0,x)$.
    The map $x \mapsto (x,0,x)$ from $\X$ to $\cY$ is the identity induction.
    
    The prequantum $1$-form $\blambda$ on $\cY$,
    uniquely determined by $\class_\omega^*(\blambda) = \Komega$,
    can be expressed on a $1$-plot $s \mapsto \gamma_s$ as:
    $\class_\omega^*(\blambda)(s \mapsto \gamma_s) = \blambda(s \mapsto \class_\omega(\gamma_s))  = \Komega(s \mapsto \gamma_s)$.
    That is,
    $\blambda(s \mapsto (x_s, t_s, x'_s))  = \Komega(s \mapsto \gamma_s)$,
    with $x_s = \gamma_s(0)$,
    $t_s = f_\omega(\gamma_s)$ and $x'_s = \gamma_s(1)$.
    But,
    \begin{align*}
      \Komega(s \mapsto \gamma_s)&  = \1^*(\alpha)(s \mapsto \gamma_s) - \0^*(\alpha)(s \mapsto \gamma_s) - [df_\omega](s \mapsto \gamma_s), \\
      & = \alpha(s \mapsto \gamma_s(1))) - \alpha(s \mapsto \gamma_s(0))) - d[s \mapsto f_\omega(\gamma_s)],\\
      & = \alpha(s \mapsto x'_s) - \alpha(s \mapsto x_s) - d[s \mapsto t_s]
    \end{align*}
    Therefore,
    $$
      \blambda(s \mapsto (x_s, t_s, x'_s)) = \alpha(s \mapsto x'_s) - \alpha(s \mapsto x_s) - d[s \mapsto t_s]
      $$
    That is,
    with reordering terms:
    $$
      \blambda = \pr_3^*(\alpha) - \pr_2^*(dt) - \pr_1^*(\alpha),
      $$
    where the $\pr_i$,
    $i = 1,2,3$,
    are the projections of $\cY \simeq \X \times \RR \times \X$ on its three factors.
    
    Since we are in a very simple case,
    we can easily check that the main properties for $(\TT_\omega, \blambda)$ are satisfied:
    
    $\ast$~\textbf{Invariance of $\blambda$ under Groupoid Composition}:
    We verify the left and right invariance of $\blambda$ using its coordinate expression:
    $y = (x, t, x') \in \cY \simeq \X \times \RR \times \X$.
    
    Let us check the left invariance.
    Let $y_0 = (x_0, t_0, x'_0) \in \cY$ and $y = (x,t,x')$.
    Then,
    $\eL(y_0)(y) = (x_0, t_0, x'_0) \cdot (x,t,x')$,
    which composes only when $x = x'_0$.
    Then $\eL(y_0)(y) = (x_0, t_0, x'_0) \cdot (x'_0,t,x') = (x_0,t_0 + t, x')$,
    that is,
    $\eL(y_0)(x'_0,t,x')= (x_0, t_0 + t, x')$.
    Now
    $\blambda \restriction {\cY_{x'_0,\ast}} = \pr_3^*(\alpha) - \pr_2^*(dt)$.
    Thus $[(x'_0,t,x') \mapsto (x_0,t_0 + t, x']^*((\pr_3^*(\alpha) - \pr_2^*(dt))) = \pr_3^*(\alpha) - \pr_2^*(dt) = \blambda \restriction {\cY_{x_0,\ast}}$.
    The same applies to right invariance.
    
    $\ast$~\textbf{Curvature of $\blambda$}:
    The differential of $\blambda$ gives $d[\pr_3^*(\alpha)] - d[\pr_2^*(dt)] - d[\pr_1^*(\alpha)] = \pr_3^*(d\alpha) - 0 - \pr_1^*(d\alpha)$,
    that is,
    $d\blambda = \pr_3^*(\omega) - \pr_1^*(\omega)$,
    which was announced by the theory.
    
    $\ast$~\textbf{Action of $\pmb{\Diff(\X, \omega)}$}:
    Let $\phi \in \Diff(\X, \omega)$,
    then $\phi^*(d\alpha) = d\alpha$ implies $d[\phi^*(\alpha)] = d\alpha$,
    and then $d[\phi^*(\alpha)-\alpha] = 0$.
    That is $\phi^*(\alpha)-\alpha$ is closed,
    but since $\X$ is simply connected $\phi^*(\alpha)-\alpha$ is exact and there exists a function $f_\phi \in \Cinfty(\X,\RR)$ such that $\phi^*(\alpha) = \alpha + df_\phi$.
    So,
    let
    $$
      \Phi[\gamma]_\omega = [\phi \circ \gamma]_\omega = \left(\phi(\gamma(0)), \int_{\phi \circ \gamma} \alpha, \phi(\gamma(1))\right).
      $$
    But $\int_{\phi \circ \gamma} \alpha = \int_\gamma \phi^*(\alpha) = \int_\gamma \alpha + df_\phi = \int_\gamma \alpha + \int_\gamma f_\phi = \int_\gamma \alpha + f_\phi(\gamma(1)) - f_\phi(\gamma(0))$.
    Therefore,
    $$
      \Phi(x,t,x') = (\phi(x), t + f_\phi(x') - f_\phi(1), \phi(x')).
      $$
    Hence,
    \begin{align*}
      \Phi^*(\blambda) &= \Phi^*(\pr_3^*(\alpha) - \pr_2^*(dt) - \pr_1^*(\alpha) \\
      & =  \Phi^*(\pr_3^*(\alpha)) - \Phi^*(\pr_2^*(dt)) - \Phi^*(\pr_1^*(\alpha)) \\
      & = (\pr_3 \circ \Phi)^*(\alpha) - (\pr_2 \circ \Phi)^*(dt) - (\pr_1 \circ \Phi)^*(\alpha)\\
      & = (\phi \circ \pr_3)^*(\alpha) - [t + f_\phi(x') - f_\phi(x)]^*(dt) -(\phi \circ \pr_1)^*(\alpha) \\
      & = \pr_3^*(\phi^*(\alpha)) - d[t + f_\phi(x') - f_\phi(x)] - \pr_1^*(\phi^*(\alpha))\\
      & = \pr_3^*(\alpha + df_\phi) - dt - d[f_\phi \circ \pr_3] + d[f_\phi \circ \pr_1] - \pr_1^*(\alpha + df_\phi) \\
      & = \pr_3^*(\alpha) -dt - \pr_1^*(\alpha) + \pr_3^*(df_\phi) - d[\pr_3^*(f_\phi)] + d[\pr_1^*(f_\phi)] - \pr_1^*(df_\phi) \\
      & = \pr_3^*(\alpha) -dt - \pr_1^*(\alpha) + \pr_3^*(df_\phi) - \pr_3^*(df_\phi) + \pr_1^*(df_\phi) - \pr_1^*(df_\phi) \\
      & = \pr_3^*(\alpha) -dt - \pr_1^*(\alpha) \\
      & = \blambda.
    \end{align*}
    This tedious computation,
    confirming that $\Diff(\X, \omega)$ acts on $\cY$ while preserving $\blambda$,
    is not without purpose.
    It shows how the function cocycle $f_\phi$ of $\Diff(\X, \omega)$,
    related to the lack of invariance of the primitive $\alpha$,
    is compensated by the variance of the isotropy such that $\Diff(\X, \omega)$ acts on $\TT_\omega$ as an automorphism of $(\TT_\omega,\blambda)$.
    
    \smallskip
    \textbf{Note.} In this exact case,
    the prequantum bundles $\cY_x$ are trivial $\RR$-bundles,
    mirroring the classical situation for exact symplectic forms on manifolds.
    This recovery of the classical structure in the simplest case demonstrates the consistency and generality of the path-space construction.
    
    The connection form $\lambda = \blambda \restriction \cY_x$ corresponds to a connection on the trivial $\RR$-bundle $\X \times \RR \to \X$ whose curvature is $\omega$.
  \end{article}
  
  %%%%%%%%%%%%%%%%%%%%%%%%%%%%%%%%%%%%%%%%%%%%%%%%%%%%%%%%%%
  %% MARK: art: Example S2
  %%%%%%%%%%%%%%%%%%%%%%%%%%%%%%%%%%%%%%%%%%%%%%%%%%%%%%%%%%
  
  \begin{article}
    \artlabel[Example: The Prequantum Groupoid for $(\S^2, \omega_{S^2})$]
    \label{art: Example S2}
    
    The $2$-sphere $\S^2$,
    equipped with its standard symplectic form $\omega_{S^2}$,
    is a fundamental example in geometric quantization.
    It is a connected and simply connected manifold,
    and $\omega_{S^2}$ is a closed $2$-form.
    The periods of $\omega_{S^2}$ are related to $\pi_2(\S^2) \simeq \ZZ$.
    For a suitable normalization (e.g., such that $\int_{\S^2} \omega_{S^2} = 2\pi$),
    the group of periods $\P_{\omega_{S^2}}$ is $2\pi\ZZ$.
    This is a discrete subgroup of $\RR$.
    Thus, $(\S^2, \omega_{S^2})$ satisfies the conditions of the main theorem.
    
    The torus of periods is $\T_{\omega_{S^2}} = \RR/\P_{\omega_{S^2}} = \RR/2\pi\ZZ \simeq \S^1$.
    
    According to the main theorem,
    the prequantum groupoid $\TT_{\omega_{S^2}}$ for $(\S^2, \omega_{S^2})$ has $\S^2$ as its objects.
    The space of morphisms is $\cY = \Paths(\S^2) / \sim_{\omega_{S^2}}$,
    where the equivalence relation $\sim_{\omega_{S^2}}$ is determined by $\omega_{S^2}$ and its periods $2\pi\ZZ$.
    The isotropy group at any point $x \in \S^2$ is isomorphic to the torus of periods: $\TT_{\omega_{S^2}, x} \simeq \T_{\omega_{S^2}} \simeq \S^1$.
    
    The space of morphisms $\cY$ has a rich structure over the product space $\S^2 \times \S^2$ via the source/target map $\ends: \cY \to \S^2 \times \S^2$.
    As discussed in Article \ref{art: Structure of the Space of Morphisms},
    this map is a diffeological fiber bundle with fiber $\T_{\omega_{S^2}} \simeq \S^1$.
    
    In this specific case,
    the prequantum line bundle over $\S^2$ (corresponding to the prequantization condition $[\omega_{S^2}] \in \HDR^2_{\mathrm{dR}}(\S^2)$ being integral) is the well-known Hopf bundle $\pi: \S^3 \to \S^2$,
    which is a principal $\S^1$-bundle.
    The space of morphisms $\cY$ of the prequantum groupoid $\TT_{\omega_{S^2}}$ can be identified with the space of isomorphisms between fibers of this principal bundle.
    This space is diffeomorphic to $(\S^3 \times \S^3) / \S^1_{\text{diag}}$,
    where $\S^1_{\text{diag}}$ is the diagonally embedded subgroup in $\S^1 \times \S^1$.
    
    The map $\ends: \cY \to \S^2 \times \S^2$ is a $\S^1$-bundle over $\S^2 \times \S^2$.
    The base space $\S^2 \times \S^2$ is a $4$-dimensional manifold,
    the fiber is $\S^1$ (a $1$-dimensional manifold),
    and the total space $\cY \simeq (\S^3 \times \S^3) / \S^1_{\text{diag}}$ is a $5$-dimensional manifold.
    The prequantum $1$-form $\blambda$ on $\cY$ is a $1$-form on this $5$-manifold whose curvature $d\blambda$ is related to $\omega_{S^2} \ominus \omega_{S^2}$ on the base $\S^2 \times \S^2$.
    
    This example provides a concrete illustration of the prequantum groupoid construction for a classical symplectic manifold,
    showing how the space of morphisms $\cY$ forms a principal $\S^1$-homogeneous-bundle over the product of the base space with itself,
    with the expected dimensions and structure related to the classical prequantum bundle.
    
    \smallskip
    \textbf{Note.} For any pair of non-antipodal points $(x,x') \in \S^2 \times \S^2$,
    the space of morphisms $\Mor_{\TT_{\omega_{S^2}}}(x,x')$,
    diffeomorphic to $\S^1$,
    can be trivialized.
    The trivialization is explicitly realized by mapping a class $[\gamma]_{\omega_{S^2}}$ (where $\gamma$ is a path from $x$ to $x'$) as follows
    $$
      [\gamma]_{\omega_{S^2}} = \bigg(x, \mathrm{e}^{i {\displaystyle \int_{\gamma_{x,x'}}^\gamma \hspace{-1em}\Komega_{S^2}}},x'\bigg),
      $$
    where $\gamma_{x,x'}$ is the unique shortest geodesic from $x$ to $x'$,
    and the integral is taken over any fixed-ends homotopy between $\gamma_{x,x'}$ and $\gamma$ in $\Paths(\S^2,x,x')$.
    This corresponds to the integral of $\omega_{S^2}$ on any surface sweeped by the homotopy.
    This identification is a local trivialization of $\cY_{\S^2}$ over the subset of non-antipodal points of the sphere.
    It connects with the physical intuition often found in the literature,
    between the abstract path-space quotient and a concrete phase value associated with paths between points.
    
  \end{article}
  
  
  %%%%%%%%%%%%%%%%%%%%%%%%%%%%%%%%%%%%%%%%%%%%%%%%%%%%%%%%%%
  %% MARK: art: Application to Symplectic Reduction
  %%%%%%%%%%%%%%%%%%%%%%%%%%%%%%%%%%%%%%%%%%%%%%%%%%%%%%%%%%
  
  \begin{article}
    \artlabel[Application to Symplectic Reduction]
    \label{art: Application to Symplectic Reduction}
    
    A central and often problematic question in geometric quantization is whether the process of quantization commutes with symplectic reduction.%
    \footnote{Reflecting on the history of this subject,
    I recall giving my first talk at Souriau's seminar more than 40 years ago, specifically on Alan Weinstein's seminal paper on symplectic reduction.
    I still vividly remember it was a Tuesday afternoon.}
    Classically,
    quantizing a symplectic manifold and then performing a reduction (Quantize then Reduce, QTR) frequently does not yield the same result as first reducing the classical system and then quantizing the reduced space (Reduce then Quantize, RTQ).
    This discrepancy is particularly acute when the reduced space is singular,
    lacking a smooth manifold structure.
    Traditional geometric quantization methods,
    heavily reliant on local Euclidean charts,
    face significant obstacles when applied directly to such singular reduced spaces.
    
    The diffeological framework,
    by providing a robust setting for differential geometry on arbitrary spaces including those with singularities,
    allows us to construct the prequantum groupoid $(\TT_\omega, \blambda)$ even for singular parasymplectic spaces arising from reduction.
    This holistic approach,
    which places the prequantum groupoid as a unified object encoding both classical and prequantum information along with the full symmetry group,
    offers a novel perspective on the reduction problem.
    This perspective resonates with the general philosophy of diffeology regarding quotient spaces,
    as discussed in \cite{GIZ25}.
    
    When considering a quotient procedure $\X/\!\sim$,
    diffeology encourages focusing on the resulting quotient space itself,
    equipped with the quotient diffeology, rather than being tied to the specific procedure that generated it.
    For instance,
    the Hopf reduction $\S^3/\S^1$ is naturally viewed as the smooth $2$-sphere $\S^2$ with its standard manifold structure (which is a particular diffeology).
    While the Hopf fibration can be retrieved by considering the set of $\S^1$ principal bundles over $\S^2$ (classified by Chern classes),
    the primary object of interest is the resulting space $\S^2$.
    
    The challenge arises when the quotient space possesses singularities and cannot be naturally equipped with a compatible manifold structure.
    In such cases,
    the resulting space might,
    at best,
    be considered a topological space, as exemplified by the quotient $\cQ_m = \CC/\ZZ_m$.
    Topologically,
    this space is homeomorphic to $\CC$.
    However,
    this topological perspective fails to capture the richer differential structure that is intuitively present,
    particularly at the singularity.
    
    This is precisely the motivation behind notions like V-manifolds (Thurston's orbifolds),
    introduced by Ichiro Satake \cite{Sat56} and shown to be equivalent to the specific class of diffeological orbifolds \cite{IKZ10}.
    Diffeology provides a framework where,
    when constructing a quotient,
    one can indeed focus on the resulting quotient diffeological space,
    regardless of the specific process.
    This applies to quotients of manifolds by compact groups,
    as shown in \cite{GIZ25},
    and is relevant to problems of symplectic reduction,
    generally denoted $\J^{-1}(0)/\G$.%
    \footnote{The map $\J$ is the moment map of a Hamiltonian equivariant action of a Lie group $\G$.
    The group $\G$ acts on the level set $\J^{-1}(0)$,
    where the restriction of the symplectic form $\omega$ is coisotropic.
    If $\J^{-1}(0)$ is a manifold such that the canonical projection to the quotient is a submersion,
    then the form $\omega \restriction \J^{-1}(0)$ descends to the quotient as a symplectic form.
    This is the Marsden-Weinstein construction \cite{MW74}.}
    
    \smallskip
    \textbf{Diffeology Principle:} A quotient $\cQ = \X/\!\sim$ is always regarded as a quotient diffeological space,
    irrespective of its construction.
    \smallskip
    
    Applying this principle to the symplectic reduction $\CC/\ZZ_m$,
    the resulting space is the cone-orbifold $\cQ_m$.
    It has been proved that $\cQ_m$ is \emph{symplectically generated} in the diffeological sense;%
    \footnote{A parasymplectic space $(\X,\omega)$ is symplectically generated if there is a generating family \cite[\textsection 1.66]{PIZ13} of plots where the pullback of $\omega$ is symplectic.}
    that is,
    there exists a closed $2$-form $\omega$ (a parasymplectic form) on $\cQ_m$ such that its pullback on $\CC$ by the projection $\pi_m : z \mapsto [z]$ is the standard symplectic form $dx \wedge dy$ \cite[\textsection 9.32]{PIZ13}.%
    \footnote{One can choose a projection $\pi_m(z) = z^m$,
    in which case $\cQ_m$ is diffeomorphic to $\CC$ equipped with the pushforward of the standard smooth diffeology on $\CC$ by this map $\pi_m$.}
    Actually,
    every parasymplectic form on $\cQ_m$ is proportional to this $\omega$ by a smooth function on $\cQ_m$.
    
    This parasymplectic diffeological space $(\cQ_m,\omega)$ is connected and simply connected;
    it is,
    in fact,
    contractible.
    The map $\rho_s: [z] \mapsto [sz]$,
    where $s \in \RR$ and $[z] \in \cQ_m$,
    is well-defined and provides a deformation retraction from $\cQ_m$ (for $s=1$) to the origin $\{[0]\}$ (for $s=0$).
    Since $\cQ_m$ is contractible,
    the parasymplectic form $\omega$ is exact \cite[\textsection 6.90]{PIZ13}.
    Let $\alpha \in \Omega^1(\cQ_m)$ be a primitive such that $\omega = d\alpha$.
    We are thus in the case previously discussed in Article~\ref{art: Exact Case}.
    The result applies directly:
    
    \begin{proposition}{}%
      ~The prequantum groupoid ${\pmb{\cQ_{m,\omega}}}$ is isomorphic to the additive groupoid
      $$
        {\pmb{\cQ_{m,\omega}}} \simeq \{(q,t,q') \in \cQ_m \times \RR \times \cQ_m \}
        \widetext{with}
        (q,t,q') \cdot (q',t',q'') = (q,t+t',q'').
        $$
      The prequantum form $\blambda$ is given by $\blambda = \pr_3^*(\alpha) - \pr_1^*(\alpha) - \pr_2^*(dt)$.
    \end{proposition}
    
    % Proposed addition after Proposition 3 in Article \ref{art: Application to Symplectic Reduction}
    This completely addresses the prequantization of the parasymplectic cone-orbifold $(\cQ_m,\omega)$ within the diffeological framework,
    thereby obviating the need to directly address the commutation of quantization and reduction procedures at the level of the base space.
    
    \textbf{Note.} A critical insight from this explicit construction is that the singular origin $0 \in \cQ_m$ does not lead to a degeneration or alteration of the isotropy group.
    For this exact case, the isotropy group $\TT_{\omega,q}$ is isomorphic to $\RR$ for all points $q \in \cQ_m$.
    This uniformity of the quantum phase information arises naturally from the inherent structure of the path-space quotient,
    as discussed in the Introduction (Note 2 of Article \ref{art: The Prequantum Groupoid}),
    demonstrating that singularities in the base space do not intrinsically affect the nature of the isotropy group itself.
    If quantum phenomena are related to such singularities, it would rather be through the representations of the group of automorphisms
    (which reveal the singular structure via the Klein stratification)
    than through the local isotropy itself.
    
  \end{article}
  
  %%%%%%%%%%%%%%%%%%%%%%%%%%%%%%%%%%%%%%%%%%%%%%%%%%%%%%%%%%
  %% MARK: art: The Uniqueness of the Quotient and the Physical Choice
  %%%%%%%%%%%%%%%%%%%%%%%%%%%%%%%%%%%%%%%%%%%%%%%%%%%%%%%%%%
  
  \begin{article}
    \artlabel[The Uniqueness of the Quotient and the Physical Choice]
    \label{art: The Physical Choice}
    
    The framework of diffeology forces a re-evaluation of the relationship between physical modeling and geometric structure.
    Consider the case of a composite system,
    such as the Deuteron,
    modeled as two indistinguishable spins.
    The classical phase space for distinguishable spins is $M = S^2 \times S^2$.
    Imposing indistinguishability leads to the quotient $\Sigma_4 = M / \fS_2$.
    
    In a context lacking the tools to handle singularities,
    one might observe that $\Sigma_4$ is homeomorphic to the complex projective plane $P^2(\CC)$.
    Since $P^2(\CC)$ carries a standard smooth manifold structure,
    it is tempting to assume that nature ``chooses'' this smooth structure by default,
    effectively smoothing out the singularity of the quotient.
    This leads to the standard quantum mechanical description of the deuteron as a single elementary particle (spin 1).
    
    However,
    in diffeology,
    this ambiguity disappears.
    The quotient $\Sigma_4$,
    equipped with the quotient diffeology,
    is a well-defined,
    rigid geometric object.
    It is distinct from $P^2(\CC)$ despite their topological equivalence;
    they are not diffeomorphic,
    and as discussed regarding the ``Boman Paradox'' \cite{PIZ25},
    topology and algebra are insufficient to define geometry.
    
    Therefore,
    we cannot transform $\Sigma_4$ into $P^2(\CC)$ by mere will or for analytical convenience.
    The choice is not between two geometries for the same system,
    but between two distinct physical hypotheses:
    \begin{enumerate}
      \item \textbf{The Elementary Hypothesis:} The system is a single,
      indivisible entity.
      The space is the homogeneous $P^2(\CC)$.
      The quantization yields an irreducible representation.
      \item \textbf{The Composite Hypothesis:} The system is composed of two identical subsystems.
      The space is the inhomogeneous $\Sigma_4$.
      The quantization yields a structure that retains the memory of the diagonal singularity (the state of coincidence).
    \end{enumerate}
    Diffeology ensures that the geometry remains faithful to the physical hypothesis.
    If one chooses the composite model,
    one must accept the singular geometry of $\Sigma_4$ and its resulting prequantum groupoid.
    The ``smoothing'' of the singularity is not a geometric regularization but a modification of the physical model itself.
  \end{article}
  
  %%%%%%%%%%%%%%%%%%%%%%%%%%%%%%%%%%%%%%%%%%%%%%%%%%%%%%%%%%
  %% MARK: art: Example Loops S3
  %%%%%%%%%%%%%%%%%%%%%%%%%%%%%%%%%%%%%%%%%%%%%%%%%%%%%%%%%%
  
  \begin{article}
    \artlabel[Example: The Prequantum Groupoid for $\pmb{\Loops(\S^3)}$]
    \label{art: Example Loops S3}
    
    The construction of the prequantum groupoid $(\TT_\omega, \blambda)$ is applicable not only to singular spaces but also to infinite-dimensional diffeological spaces,
    provided they satisfy the conditions of the main theorem (connected, simply connected, with a closed $2$-form having discrete periods).
    A relevant example is the loop space of the three-sphere,
    $\X = \Loops(\S^3)$.
    
    The space $\X = \Loops(\S^3)$ is an infinite-dimensional diffeological space.
    Its path-connectedness is given by $\pi_0(\Loops(\S^3)) \simeq \pi_1(\S^3) = \{0\}$;
    thus $\X$ is connected.
    Its simple connectedness is given by $\pi_1(\Loops(\S^3)) \simeq \pi_2(\S^3) = \{0\}$;
    thus $\X$ is simply connected.
    
    Consider the volume $3$-form $\vol$ on $\S^3$.
    Let $\omega = \CHO\vol \in \Omega^2(\X)$ be the $2$-form on $\X$ obtained by applying the chain-homotopy operator.
    The $2$-form $\omega$ is closed,
    which follows from $d\vol = 0$ and the properties of the chain-homotopy operator.
    The group of periods $\P_\omega$ of the form $\omega$ is the group of periods of the $1$-form $\Komega \restriction \Loops(\Loops(\X))$.
    This group is a homomorphic image of $\pi_1(\Loops(\X)) \simeq \pi_2(\X) \simeq \pi_3(\S^3) = \ZZ$.
    The periods of $\omega$ are related to integrals of $\vol$ over cycles in $\pi_3(\S^3)$.
    For a suitable normalization of the volume form $\vol$ on $\S^3$ (e.g., such that $\int_{\S^3} \vol = 1$),
    the group of periods $\P_\omega$ is isomorphic to $\ZZ$.
    
    The torus of periods is $\T_\omega = \RR/\P_\omega = \RR/\ZZ \simeq \S^1$.
    
    According to the main theorem,
    the prequantum groupoid $\TT_\omega$ for $(\X, \omega)$ has $\X = \Loops(\S^3)$ as its objects.
    The space of morphisms is $\cY = \Paths(\X) / \sim_\omega = \Paths(\Loops(\S^3)) / \sim_\omega$,
    where the equivalence relation $\sim_\omega$ is determined by the $2$-form $\omega$ and its periods $\P_\omega = \ZZ$.
    
    The isotropy group at any object $\ell \in \X$ (i.e., at any loop $\ell \in \Loops(\S^3)$) is isomorphic to the torus of periods:
    $\TT_{\omega, \ell} \simeq \T_\omega \simeq \S^1$.
    This means that at each point in the infinite-dimensional space $\Loops(\S^3)$,
    the quantum phase information is captured by an $\S^1$ group.
    
    The prequantum $1$-form $\blambda$ on $\cY$ is the unique form such that $\class_\omega^*(\blambda) = \Komega$.
    
    Explicitly representing the space of morphisms $\cY = \Paths(\Loops(\S^3)) / \sim_\omega$ in a simpler form (like a product space) is generally not possible due to the infinite-dimensional nature of $\X$ and the non-exactness of $\omega$ (since $\P_\omega = \ZZ \neq \{0\}$).
    The space $\cY$ is genuinely the quotient space defined by the equivalence relation.
    
    This example illustrates how the prequantum groupoid construction extends to infinite-dimensional settings.
    The situation shares some analogy with the prequantization of the $2$-sphere $(\S^2, \omega_{S^2})$,
    where $\S^2$ is connected and simply connected,
    and the periods of the standard symplectic form are related to $\pi_2(\S^2) \simeq \ZZ$,
    leading to an $\S^1$ isotropy group.
    The key difference here is the infinite dimensionality of the base space $\X = \Loops(\S^3)$.
    
  \end{article}
  
  %%%%%%%%%%%%%%%%%%%%%%%%%%%%%%%%%%%%%%%%%%%%%%%%%%%%%%%%%%
  %% MARK: art: The Moment Map on the Prequantum Groupoid
  %%%%%%%%%%%%%%%%%%%%%%%%%%%%%%%%%%%%%%%%%%%%%%%%%%%%%%%%%%
  
  \begin{article}\artlabel[The Moment Map on the Prequantum Groupoid]
    \label{art:Momentum Map}
    
    Now that we have this symmetry group $\Aut(\TT_\omega, \blambda)$,
    we can define a moment map using its action,
    and here are its properties and what it tells us.
    
    The prequantum groupoid $(\TT_\omega, \blambda)$ provides a geometric framework to study the system $(\X, \omega)$.
    As established in Article \ref{art:Automorphisms Induced by Symmetries},
    its group of automorphisms $\G_\omega = \Aut(\TT_\omega, \blambda)$ is diffeologically isomorphic to the group of $\omega$-preserving diffeomorphisms of $\X$, $\Diff(\X, \omega)$.
    This group $\G_\omega$ acts smoothly on the space of morphisms $\cY = \Mor(\TT_\omega)$,
    preserving the prequantum $1$-form $\blambda$.
    
    According to the general theory of moment maps in diffeology developed in \cite{PIZ10} (see also \cite[Chap.\ 9]{PIZ13}),
    the existence of a $\G_\omega$-invariant $1$-form $\blambda$ on $\cY$ ensures the existence of a moment map $\Psi_\omega : \cY \to \cG_\omega^*$,
    where $\cG_\omega^*$ is the space of left-invariant $1$-forms on the diffeological group $\G_\omega$ (the space of momenta \cite[Art.~2.6]{PIZ10}).
    This map is defined for $y \in \cY$ by $\Psi_\omega(y) = \hat y^*(\blambda)$,
    where $\hat y : \G_\omega \to \cY$ is the orbit map $\hat y(\Phi) = \Phi(y)$.
    
    This moment map on the space of morphisms $\cY$ is directly related to the paths moment map defined in \cite[Art.~3.1]{PIZ10} for the action of $\G_\omega$ on $\X$.
    Let $\Psi_{\text{paths}} : \Paths(\X) \to \cG_\omega^*$ be the paths moment map given by $\Psi_{\text{paths}}(\gamma) = \hat\gamma^*(\Komega)$,
    where $\hat\gamma(\phi) = \phi \circ \gamma$.
    The map $\Psi_\omega$ is precisely the projection of $\Psi_{\text{paths}}$ onto the quotient $\cY$:
    
    \begin{proposition}{}
      The moment map $\Psi_\omega : \cY \to \cG_\omega^*$ satisfies $\Psi_\omega \circ \class_\omega = \Psi_{\text{paths}}$.
    \end{proposition}
    
    \begin{proof}
      Let $\gamma \in \Paths(\X)$,
      and let $y = \class_\omega(\gamma) \in \cY$.
      By definition of $\Psi_\omega$, $\Psi_\omega(y) = \hat y^*(\blambda)$,
      where $\hat y : \G_\omega \to \cY$ is the orbit map $\hat y(\Phi) = \Phi(y)$ for $\Phi \in \G_\omega$.
      Recalling that $\G_\omega = \Aut(\TT_\omega, \blambda)$ is isomorphic to $\Diff(\X,\omega)$ via $\Phi \mapsto \phi$ such that $\Phi([\gamma']_\omega) = [\phi \circ \gamma']_\omega$,
      the orbit map $\hat y$ is given by $\hat y(\Phi) = [\phi \circ \gamma]_\omega = \class_\omega(\phi \circ \gamma)$.
      Thus,
      $\hat y = \class_\omega \circ \hat\gamma$,
      where $\hat\gamma : \G_\omega \to \Paths(\X)$ is the orbit map $\hat\gamma(\phi) = \phi \circ \gamma$ as defined for $\Psi_{\text{paths}}$.
      Now we compute the pullback:
      $$
        \Psi_\omega(\class_\omega(\gamma)) = \hat y^*(\blambda) = (\class_\omega \circ \hat\gamma)^*(\blambda) = \hat\gamma^*(\class_\omega^*(\blambda)).
        $$
      By definition of $\blambda$ (Theorem \ref{sec:Main Theorem: The Prequantum Groupoid}, part 3),
      $\class_\omega^*(\blambda) = \Komega$.
      Substituting this gives:
      $$
        \Psi_\omega(\class_\omega(\gamma)) = \hat\gamma^*(\Komega).
        $$
      The right side is precisely $\Psi_{\text{paths}}(\gamma)$.
      Therefore, $\Psi_\omega \circ \class_\omega = \Psi_{\text{paths}}$.
    \end{proof}
    \smallskip
    
    The paths moment map $\Psi_{\text{paths}}$ is an additive homomorphism from the path concatenation to $(\cG_\omega^*, +)$ \cite[Art.~4.3]{PIZ10}.
    Since $\Psi_\omega$ factors through $\class_\omega$ and $\class_\omega$ preserves concatenation,
    $\Psi_\omega$ is an additive homomorphism from the groupoid $\TT_\omega$ to $(\cG_\omega^*, +)$:
    $$
      \Psi_\omega(y \cdot y') = \Psi_\omega(y) + \Psi_\omega(y')
      $$
    for any composable $y, y' \in \cY$.
    
    The additive property $\Psi_\omega(y \cdot y') = \Psi_\omega(y) + \Psi_\omega(y')$ implies that the value of $\Psi_\omega(y)$ for $y \in \cY$ depends only on the endpoints of $y$.
    This means $\Psi_\omega$ descends to a map $\psi_\omega : \X \times \X \to \cG_\omega^*$.
    This descent occurs if and only if $\Psi_\omega$ is constant on the equivalence classes $[\gamma]_\omega$ corresponding to loops,
    which is equivalent to $\Psi_\omega(\ell)=0$ for all loops $\ell \in \Loops(\X)$ (viewed as elements of $\cY$ via $\class_\omega$).
    The set of values $\Psi_{\text{paths}}(\ell) = \Psi_\omega(\class_\omega(\ell))$ for $\ell \in \Loops(\X)$ is the universal holonomy group $\Gamma_\omega = \Psi_{\text{paths}}(\Loops(\X))$ for the action of $\G_\omega$ on $\X$ \cite[Art.~3.7]{PIZ10}.
    
    Since $\X$ is connected and simply connected,
    its fundamental group $\pi_1(\X)$ is trivial.
    The universal holonomy group $\Gamma_\omega$ is an additive subgroup of $\cG_\omega^*$ which is a homomorphic image of $\pi_0(\Loops(\X))$ \cite[Art.~3.7]{PIZ10}.
    Since $\X$ is simply connected and connected,
    $\pi_0(\Loops(\X))$ is trivial.
    Therefore,
    for a simply connected $\X$,
    the universal holonomy $\Gamma_\omega$ is trivial,
    $\Gamma_\omega = \{0\}$.
    This property will not be preserved in the general case of $\pi_1(\X) \neq \{0\}$.
    In the general case,
    the universal holonomy $\Gamma_\omega$ will typically be non-trivial.
    
    The fact that $\Gamma_\omega=\{0\}$ for simply connected $\X$ implies that $\Psi_\omega(\ell)=0$ for all loops $\ell \in \Loops(\X)$,
    and thus $\Psi_\omega : \cY \to \cG_\omega^*$ descends to a well-defined map on $\X \times \X$.
    This descended map is the \emph{two-points moment map}:
    $$
      \psi_\omega : \X \times \X \to \cG_\omega^*
      \quad \text{defined by} \quad
      \psi_\omega(x,x') = \Psi_\omega(y),
      $$
    for any $y \in \cY$ with $\ends(y) = (x,x')$.
    This $\psi_\omega$ is the universal two-points moment map associated with $\omega$ (compare \cite[Art.~4.1]{PIZ10}).
    It satisfies the Chasles cocycle relation $\psi_\omega(x,x'') = \psi_\omega(x,x') + \psi_\omega(x',x'')$ for $x,x',x'' \in \X$.%
    \footnote{~This moment map $\psi_\omega$ can be interpreted as the moment map for the diffeological space $\X \times \X$ equipped with the $2$-form $\omega \ominus \omega = \pr_1^*\omega - \pr_2^*\omega$ and the diagonal action of $\G_\omega$.}
    
    From the two-points moment map $\psi_\omega$,
    we can derive the \emph{one-point moment map} $\mu_\omega : \X \to \cG_\omega^*$ by choosing a base point $x_0 \in \X$.
    As shown in \cite[Art.~5.1]{PIZ10},
    $\mu_\omega$ is a primitive of $\psi_\omega$,
    unique up to an additive constant $\varepsilon \in \cG_\omega^*$,
    satisfying $\psi_\omega(x,x') = \mu_\omega(x') - \mu_\omega(x)$.
    A standard choice is $\mu_\omega(x) = \psi_\omega(x_0,x)$.
    
    The equivariance properties of these moment maps and the associated Souriau's cocycle $\theta_\omega \in \Cinfty(\G_\omega, \cG_\omega^*)$ are described in detail in \cite[Art.~5.2]{PIZ10}.
    For simply connected $\X$,
    the holonomy $\Gamma_\omega$ is trivial,
    which means the target space for Souriau's cocycle and moment maps is $\cG_\omega^*$ (as opposed to $\cG_\omega^*/\Gamma_\omega$ in the general case).
    The universal Souriau class $\sigma_\omega$,
    which lives in $\H^1(\G_\omega, \cG_\omega^*/\Gamma_\omega)$,
    is therefore an element of $\H^1(\G_\omega, \cG_\omega^*)$.
    The triviality of this cohomology class implies the existence of an equivariant moment map $\mu_\omega$ \emph{op.~cit.}
    However,
    the triviality of $\Gamma_\omega$ does not in general imply the triviality of the entire cohomology group $\H^1(\G_\omega, \cG_\omega^*)$.
    Thus,
    for a simply connected $\X$,
    while the holonomy is trivial,
    Souriau's class $\sigma_\omega$ may be non-trivial,
    corresponding to non-trivial central extensions in the classical picture.
    
  \end{article}
  
  %%%%%%%%%%%%%%%%%%%%%%%%%%%%%%%%%%%%%%%%%%%%%%%%%%%%%%%%%%
  %% MARK: art: The Geometric Nature of Quantum States
  %%%%%%%%%%%%%%%%%%%%%%%%%%%%%%%%%%%%%%%%%%%%%%%%%%%%%%%%%%
  
  \begin{article}\artlabel[The Geometric Nature of Quantum States]
    \label{art:The Geometric Nature of Quantum States}
    
    The construction of the prequantum groupoid $\TT_\omega$ suggests a shift in the definition of a quantum state.
    Rather than starting with a Hilbert space of sections,
    we look for objects that are intrinsic to the groupoid structure itself.
    
    A natural candidate arises from the algebraic structure of the morphisms.
    Let $\V$ be a complex vector space (typically $\CC$).
    A smooth function $\Psi: \cY \to \V$ is called \emph{multiplicative} if,
    for any composable pair $y_1, y_2 \in \cY$,
    we have:
    $$
      \Psi(y_1 \cdot y_2) = \Psi(y_1) \cdot \Psi(y_2).
      $$
    This definition directly encodes the quantum phase.
    For any loop $\tau$ in the isotropy group $\TT_{\omega, x} \simeq \T_\omega$,
    the condition implies $\Psi(\tau \cdot y) = \chi(\tau)\Psi(y)$,
    where $\chi$ is a character of the torus of periods.
    
    These multiplicative functions are the geometric "atoms" of the quantum theory.
    They correspond to the "pure states" or "semiclassical amplitudes" associated with the geometry.
    While the full analytical machinery (superposition,
    operators) requires the construction of a convolution algebra (as discussed in Part II),
    these functions constitute the fundamental geometric dual of the prequantum groupoid.
    They demonstrate that the quantum state is not an external addition to the theory,
    but a covariant representation of the space of motions itself.
  \end{article}
  
  %%%%%%%%%%%%%%%%%%%%%%%%%%%%%%%%%%%%%%%%%%%%%%%%%%%%%%%%%%
  %% MARK: art: From Geometry to Algebra
  %%%%%%%%%%%%%%%%%%%%%%%%%%%%%%%%%%%%%%%%%%%%%%%%%%%%%%%%%%
  \begin{article}\artlabel[From Geometry to Algebra]
    \label{art:From Geometry to Algebra}
    
    The prequantum groupoid $\TT_\omega$ constitutes the objective geometric reality of the quantum system.
    However,
    to connect this geometry with the analytical machinery of quantum mechanics (operators,
    spectra,
    states),
    we must transition from the \emph{points} of the space to the \emph{functions} defined upon it.
    This transition occurs in two stages:
    the identification of elementary amplitudes and the construction of the algebra of observables.
    
    1.~\emph{Elementary Amplitudes (Multiplicative Functions).}
    A natural class of objects arises directly from the groupoid structure.
    A smooth,
    complex-valued function $\Psi$ on the space of morphisms $\cY$ is called \emph{multiplicative} if,
    for any composable pair $y, y' \in \cY$,
    we have:
    $$
      \Psi(y \cdot y') = \Psi(y) \cdot \Psi(y').
      $$
    This condition enforces that $\Psi$ transforms according to a character of the isotropy group $\T_\omega$.
    Specifically,
    for any loop $\tau \in \TT_{\omega, \src(y)} \simeq \T_\omega$,
    we have $\Psi(\tau \cdot y) = \chi(\tau)\Psi(y)$.
    These functions represent the "pure phases" or "semiclassical states" of the system.
    They are the geometric "atoms" of the theory,
    corresponding to flat sections in the bundle picture.
    
    2.~\emph{The Prequantum Algebra.}
    Physical states and observables,
    however,
    rely on the principle of linear superposition.
    They are not single paths,
    but weighted sums (integrals) over the space of motions.
    We define the \emph{Prequantum Algebra} $\cA_\omega$ as the space of compactly supported complex densities on $\cY$.
    The groupoid structure induces a natural \emph{convolution product} on this space:
    $$
      (f * g)(y) = \int_{y_1 \cdot y_2 = y} f(y_1) g(y_2).
      $$
    This operation transforms the geometry of the groupoid into a non-commutative $*$-algebra.
    The multiplicative functions defined above play a distinguished role here:
    they are the elementary characters (or "plane waves") of this algebra.
    The general elements of the algebra are linear combinations of these geometric phases.
    
    3.~\emph{The Threshold of Analysis.}
    The algebra $\cA_\omega$ retains the full covariance of the classical system.
    The group of symmetries $\Diff(X, \omega)$ acts by automorphisms on the groupoid,
    and therefore by automorphisms on the algebra.
    
    The transition to standard quantum mechanics requires a final step:
    \emph{Polarization}.
    This consists of selecting a specific representation of this algebra on a Hilbert space.
    This step is analytically necessary but geometrically traumatic,
    as it often breaks the symmetry of the groupoid.
    By stopping at the algebra $\cA_\omega$,
    we preserve the purity of the geometric signal.
    We have constructed the universal container for the physics,
    standing exactly at the threshold where geometry yields to analysis.
  \end{article}
  
  %%%%%%%%%%%%%%%%%%%%%%%%%%%%%%%%%%%%%%%%%%%%%%%%%%%%%%%%%%
  %% MARK: art: Relation to Classical Central Extensions
  %%%%%%%%%%%%%%%%%%%%%%%%%%%%%%%%%%%%%%%%%%%%%%%%%%%%%%%%%%
  
  \begin{article}
    \artlabel[Relation to Classical Central Extensions]
    \label{art:Automorphism Group Extension} % Keep the label for potential internal references
    
    It is natural to question how the isomorphism $\Aut(\TT_\omega, \blambda) \simeq \Diff(\X, \omega)$ relates to the well-known phenomenon of non-trivial central extensions arising in the classical geometric quantization of symmetries, even on simply connected symplectic manifolds.
    For instance,
    the prequantization of the standard symplectic structure on $\RR^{2n}$ leads to the Heisenberg group extension for the translation group,
    and the prequantization of the sphere $\S^2$ leads to the spin group $\mathrm{Spin}(3) \simeq \mathrm{SU}(2)$ as a central extension of $\SO(3)$.
    
    The key to understanding this lies in distinguishing the object to which the symmetries are lifted:
    
    $\ast$~In the \textbf{classical prequantum bundle picture},
    central extensions emerge when considering the lift of a symmetry group (or its Lie algebra) acting on the base manifold $\X$ to the total space of the \emph{prequantum principal $\U(1)$-bundle} $P \to \X$,
    or equivalently, to the space of sections of the associated line bundle. The failure of the lifted Lie bracket to match the lift of the base Lie bracket results in a cocycle valued in the Lie algebra of the fiber group ($\mathfrak{u}(1)$),
    leading to a central extension of the symmetry Lie algebra and group by $\U(1)$.
    
    $\ast$~In the \textbf{prequantum groupoid picture},
    the isomorphism $\Aut(\TT_\omega, \blambda) \simeq \Diff(\X, \omega)$ states that the full group of $\omega$-preserving diffeomorphisms $\phi$ lifts \emph{isomorphically} to act as automorphisms $\Phi$ of the \emph{entire prequantum groupoid} $(\TT_\omega, \blambda)$ that preserve the prequantum 1-form $\blambda$.
    The object being acted upon without central extension here is the groupoid's space of morphisms $\cY = \Paths(\X)/\!\sim_\omega$ endowed with the form $\blambda$.
    
    The classical prequantum bundle $\Y \to \X$
    (or, in our framework, the generalized prequantum bundle $\cY_x = \Mor(\TT_\omega, x, \ast) \to \X$, obtained by fixing a source point $x$)
    is \emph{not} the same object as the full groupoid $\TT_\omega$.
    The groupoid $\TT_\omega$ encompasses all such principal bundles $\cY_x$ for varying $x$,
    and the arrows of $\TT_\omega$ can be seen as isomorphisms between fibers of these bundles over different points.
    
    The central extensions observed in the classical setting \emph{will reappear} when one descends from the groupoid level to consider actions on structures derived from the groupoid over the base space $\X$.
    For example,
    when considering the lifting of symmetries to the total space of the principal bundle $\cY_x \to \X$ or,
    crucially,
    to the space of sections of an associated bundle over $\X$ (which would be the candidate quantum Hilbert space),
    the central extension structure related to the period group $\T_\omega$ will indeed typically emerge.
    
    Therefore,
    the result $\Aut(\TT_\omega, \blambda) \simeq \Diff(\X, \omega)$ does not contradict classical findings.
    Instead,
    it provides a potentially deeper insight: the groupoid $(\TT_\omega, \blambda)$ serves as a fundamental,
    extension-free prequantum object where the entire symmetry group $\Diff(\X, \omega)$ is faithfully represented.
    The "quantic anomaly" manifested by central extensions in the classical picture is then understood as a phenomenon that arises when one projects or realizes this fundamental structure and its symmetries within the context of bundles or representation spaces built over the base space $\X$.
    The essential quantic information,
    captured by the period group $\T_\omega$, resides fundamentally in the isotropy group of the groupoid,
    waiting to express itself as an extension upon descent to structures defined over $\X$. This perspective offers significant conceptual clarity,
    especially when dealing with singular spaces where classical bundle constructions are challenging.
    
    \smallskip
    \textbf{Note.} This phenomenon of symmetry breaking is reminiscent of passing from an equivariant two-point moment map $\psi(x,x')$ to a no longer equivariant one-point moment map $\mu(x) = \psi(x_0,x) + c$,
    which is the general solution of the equation $\psi(x,x') = \mu(x') - \mu(x)$.
    The one-point moment map is no longer equivariant under the linear coadjoint action but with the affine coadjoint action modified by the cocycle $\theta(g) = \psi(x_0,g(x_0)) + \Delta c$.
  \end{article}
  
  
  %%%%%%%%%%%%%%%%%%%%%%%%%%%%%%%%%%%%%%%%%%%%%%%%%%%%%%%%%%
  %% MARK: art: Future Directions
  %%%%%%%%%%%%%%%%%%%%%%%%%%%%%%%%%%%%%%%%%%%%%%%%%%%%%%%%%%
  
  \begin{article}
    \artlabel[Future Directions]
    \label{art:Future Directions}
    
    This paper has established the construction of the prequantum groupoid for connected and simply connected diffeological spaces.
    This constitutes the first step in a broader program to reformulate geometric quantization within the framework of diffeology.
    
    The extension of this construction to the general case (non-simply connected spaces),
    as well as the transition to the algebraic framework,
    are the subject of the sequel to this work:
    \emph{Geometric Quantization by Paths, Part II: The General Case}.
    In that work,
    we address:
    \begin{enumerate}
      \item \textbf{Homotopy and Periods:} We show how the fundamental group and the non-trivial homotopy of the space generate "surfacic periods."
      By incorporating these into the definition of the \emph{Total Group of Periods},
      we construct a prequantum groupoid for arbitrary connected spaces,
      resolving the obstructions associated with non-trivial topology.
      \item \textbf{The Prequantum Algebra:} We transition from the geometry of the groupoid to the analysis of the system by constructing the \emph{Convolution Algebra} of the groupoid.
      This algebra constitutes the intrinsic observable algebra of the system,
      retaining the full covariance of the classical symmetries.
    \end{enumerate}
    
    Beyond these immediate results,
    the groupoid framework opens several avenues for long-term research:
    
    $\ast$~\textbf{The Feynman Connection.}
    The path-based nature of our construction suggests that the prequantum groupoid is the "geometric crystallization" of the Feynman path integral.
    Formalizing the measure-theoretic link between the groupoid structure and the path integral amplitude remains a major open question.
    
    $\ast$~\textbf{Polarization and the Hilbert Space.}
    To recover standard quantum mechanics,
    one must select a specific representation of the prequantum algebra.
    This requires adapting the theory of \emph{polarizations} (Lagrangian or complex) to the groupoid setting,
    interpreting them as the choice of a "screen" on which the quantum geometry is projected.
    
    $\ast$~\textbf{The Maslov Correction.}
    In the classical bundle approach,
    the metaplectic correction is required to obtain the correct spectrum (e.g., for the harmonic oscillator).
    Investigating how this topological correction manifests intrinsically within the space of paths and the structure of the groupoid is essential.
    
    $\ast$~\textbf{Coadjoint Orbits.}
    Applying this construction to the coadjoint orbits of infinite-dimensional Lie groups (such as the Virasoro group or groups of diffeomorphisms) could provide new insights into their quantization,
    where standard manifold techniques often fail.
    
  \end{article}
  
  %%%%%%%%%%%%%%%%%%%%%%%%%%%%%%%%%%%%%%%%%%%%%%%%%%%%%%%%%%
  %% MARK: Bibliography
  %%%%%%%%%%%%%%%%%%%%%%%%%%%%%%%%%%%%%%%%%%%%%%%%%%%%%%%%%%
  
  \begin{thebibliography}{IKZ10} % Use 99 or a number larger than your total entries
    
    \bibitem[DI83]{DI83}
    Paul Donato and Patrick Iglesias.
    {\em Exemple de groupes diff\'erentiels : flots irrationnels sur le tore}.
    Preprint CPT-83/P.1524, Centre de Physique Th\'eorique, Marseille, July 1983. Published in {\em Comptes Rendus de l'Acad\'emie des Sciences}, 301(4), Paris, 1985.
    \texttt{http://math.huji.ac.il/~piz/documents/EDGDFISLT.pdf}
    
    \bibitem[GIZ23]{GIZ23}
    Serap G{\"u}rer and Patrick Iglesias-Zemmour.
    {\em Orbifolds as stratified diffeologies}.
    {Differential Geometry and its Applications}, 86:101969, 2023.
    
    \bibitem[GIZ25]{GIZ25}
    Serap Gürer and Patrick Iglesias-Zemmour,
    \emph{On Diffeology of Orbit Spaces},
    (Preprint available at \texttt{http://math.huji.ac.il//~piz/documents/ODOOS.pdf}).
    
    \bibitem[Hor24]{Hor24}
    Peter A. Horvathy.
    {\em Prequantisation from the path integral viewpoint}.
    Preprint arXiv:2402.17629, 2024.
    
    \bibitem[PIZ85]{PIZ85}
    Patrick Iglesias.
    {\em Fibr\'es diff\'eologiques et homotopie}.
    Th\`ese de doctorat d'\'etat, Universit\'e de Provence, Marseille, 1985.
    
    \bibitem[PIZ95]{PIZ95}
    Patrick Iglesias.
    {\em La trilogie du moment}.
    Ann. Inst. Fourier, 45(3):825--857, 1995.
    
    \bibitem[IKZ10]{IKZ10}
    Patrick Iglesias, Yael Karshon, and Moshe Zadka.
    {\em Orbifolds as diffeologies}.
    Transactions of the American Mathematical Society, Vol. 362, Number 6, Pages 2811--2831, Providence R.I., 2010.
    
    \bibitem[PIZ10]{PIZ10}
    Patrick Iglesias-Zemmour,
    \emph{The moment maps in diffeology},
    Mem. Amer. Math. Soc. \textbf{204} (2010), no. 962, viii+101 pp.
    
    \bibitem[PIZ13]{PIZ13}
    Patrick Iglesias-Zemmour.
    {\em Diffeology}.
    Mathematical Surveys and Monographs, 185, Am. Math. Soc., Providence, RI, 2013.
    (Revised version by Beijing World Publishing Corp., Beijing, 2022.)
    
    \bibitem[PIZ22]{PIZ22}
    Patrick Iglesias-Zemmour.
    {\em Diffeology}.
    Reprint by Beijing World Publishing Corporation, Ltd. Beijing Branch, 2022. ISBN 978-7-5192-9608-7.
    Accessible at \\
    \url{http://math.huji.ac.il/~piz/documents/Diffeology.pdf}
    
    \bibitem[PIZ15]{PIZ15}
    \bysame
    {\em Example of singular reduction in symplectic diffeology}.
    Proceedings of the American Mathematical Society, 144(3):1309--1324, 2016.
    
    \bibitem[IZP21]{IZP21}
    Patrick Iglesias-Zemmour and Elisa Prato.
    {\em Quasifolds, diffeology and noncommutative geometry}.
    J. Noncommut. Geom., 15(2):735--759, 2021.
    
    \bibitem[PIZ24]{PIZ24}
    Patrick Iglesias-Zemmour,
    \emph{{\v C}ech-de Rham bicomplex in diffeology},
    Israel J. Math. \textbf{259} (2024), no. 1, 239--276.
    
    \bibitem[PIZ25]{PIZ25}
    \bysame,
    \emph{The Boman Paradox --- When topology plus algebra do not define geometry},
    To appear in the Mathematical Intelligencer. Preprint 2025. Available at \url{http://math.huji.ac.il/~piz/documents/TBP.pdf}.
    
    \bibitem[MW74]{MW74}
    Jerrold E. Marsden and Alan Weinstein,
    \emph{Reduction of symplectic manifolds with symmetry},
    Rep. Mathematical Phys. \textbf{5} (1974), no. 1, 121--130.
    
    \bibitem[Sat56]{Sat56}
    Ichiro Satake,
    \emph{On a generalization of the notion of manifold},
    Proceedings of the National Academy of Sciences \textbf{42} (1956), 359--363.
    
    \bibitem[Sou70]{Sou70}
    Jean-Marie Souriau.
    {\em Structure des syst\`emes dynamiques}.
    Dunod, Paris, 1970.
    
    \bibitem[Sou75]{Sou75}
    \bysame %Jean-Marie Souriau.
    \newblock \emph{Construction explicite de l'indice de Maslov. Applications}.
    \newblock In: Group theoretical methods in physics (Fourth Internat. Colloq., Nijmegen, 1975), Lecture Notes in Phys., Vol. 50, Springer, Berlin, 1976, pp. 117--148.
    
    \bibitem[Sou04]{Sou04}
    Jean-Marie Souriau,
    \emph{C’est Quantique~? donc c’est géométrique},
    in: Dominique Flament (dir),
    \emph{Feuilletages -quantification géométrique : textes des journées d’étude des 16 et 17 octobre 2003, Paris, Fondation Maison des Sciences de l’Homme},
    Série Documents de travail (Équipe F2DS), 2004.
    
  \end{thebibliography}
  
  %%%%%%%%%%%%%%%%%%%%%%%%%%%%%%%%%%%%%%%%%%%%%%%%%%%%%%%%%%
  %%
  %% MARK: End Private comments
  %%
  %%%%%%%%%%%%%%%%%%%%%%%%%%%%%%%%%%%%%%%%%%%%%%%%%%%%%%%%%%
  
  
\end{document}
